\documentclass[11pt]{article}
\usepackage[a4paper,margin=25mm]{geometry}
\usepackage{amsmath,amssymb,amsthm,booktabs}
\usepackage[british]{babel}
\usepackage{url}
\usepackage{xcolor,graphicx}
\usepackage[round,authoryear]{natbib}
\setlength{\bibhang}{1.5em}
\setlength{\bibsep}{2pt}
\renewcommand{\bibfont}{\small}
\usepackage[hidelinks]{hyperref}
\hypersetup{pdfauthor={Amanjeet Singh},
pdftitle={Certified curvature corrections for quantum basket pricing},
pdfsubject={Sharp error envelopes and cancellation-aware resource comparisons},
pdfkeywords={quantum finance, basket pricing, Gaussian conditioning}}
\setlength{\emergencystretch}{3em}
\newtheorem{theorem}{Theorem}[section]
\newtheorem{lemma}[theorem]{Lemma}
\newtheorem{proposition}[theorem]{Proposition}
\newtheorem{corollary}[theorem]{Corollary}
\newcommand{\E}{\mathbb E}
\newcommand{\R}{\mathbb R}
\newcommand{\ph}{\varphi}
\newcommand{\ind}{\mathbf 1}
\DeclareMathOperator{\Var}{Var}
\title{Certified curvature corrections for quantum basket pricing\\[3mm]
\large Sharp error envelopes and cancellation-aware resource comparisons}
\author{Amanjeet Singh}
\date{28 October 2026}
\begin{document}
\maketitle
\begin{abstract}
\noindent
Quantum resource bounds for capped basket pricing depend on how the payoff
is represented coherently. Under a classical no-arbitrage lognormal model,
I combine Gaussian conditioning, payoff-boundary curvature and moment bounds
to certify pricing errors and compare coherent preparations under common
information, accuracy and confidence requirements. Complete logical T-gate
accounting includes state preparation, reversible arithmetic and reflections,
with the same cancellation-aware optimisation applied to every eligible rival.
Across four stressed synthetic contracts, shared-threshold preparation reduces
the resulting upper bounds by 2.34 to 10.00 per cent relative to the cheaper
equally optimised implemented rival. Matched Chebyshev diagnostics instead
favour numerical conditioning in mean resource charge on the stipulated
probability grid. Classical methods certify all 40 prices in the broader
test set. These reductions concern a declared finite architecture and
precision pool and do not establish practical quantum advantage.
\end{abstract}

\smallskip
\noindent\begin{minipage}{\linewidth}
\small\emergencystretch=2em
\textit{Keywords:} financial mathematics; basket option; Gaussian conditioning;
payoff boundary curvature; Gaussian transport; reversible arithmetic;
amplitude estimation; quantum finance.\par
\smallskip
\textit{JEL classification:} C63, C65, G13.\par
\end{minipage}

\section{Research question and relation to the literature}
\label{sec:question}

I ask whether coherent representations of controlled basket-price
corrections reduce complete quantum pricing cost when strong opponents
receive the same financial model, analytic controls, derivative constants,
moment information, accuracy and confidence. I compare an integrated
boundary kernel, a stable two-node quadrature kernel, raw direct second
correction, numerical positive first correction, a shared threshold
first correction and a common-normal capped-payoff first correction. Quantum
mechanics enters through coherent preparation, reversible arithmetic,
interference and measurement. The underlying risk-neutral law remains
classical, and the priced claim has an explicit no-arbitrage interpretation.

Geometric conditioning and financial smoothing are established
\citep{curran,dingechormann,bayer}. Polynomial and conditional Taylor pricing
have close antecedents \citep{ju,olivaresalvarez2016}; general basket bounds
also supply substantive classical competitors \citep{caldana}. Curran's geometric-tail evaluation and approximate below-strike gap
calculation already establish the conditioning architecture.
\footnote{I consulted the complete final Curran scan and the final publisher
analysis of Olivares--Alvarez, supplemented by its distinct author version
where the source ledger identifies formula gaps. Ju's complete article
remains unavailable and is not a completed source-level priority exclusion.} The falling log-price derivative operator is
already given by \citet{carr2001}. Neither Gaussian call pricing, this
operator nor Gaussian quadrature is a new generic mechanism.

Amplitude estimation and amplification provide the established estimation
framework \citep{brassard,montanaro}. Published quantum derivative-pricing
analyses cover signed or unbounded payoffs and substantial implementation
costs \citep{rebentrost,chakrabarti}. In particular,
\citet{chenlineufeld2026} analyses piecewise-affine payoffs of linear growth,
state preparation, financial grids and tails. \footnote{I verified its published \emph{Mathematical Finance} record;
the detailed sections actually consulted are from the authors' May 2026
preprint, as recorded in the source ledger.}
Its state-preparation assumptions and elementary-gate complexity are not
identical to the realised logical-T contract here. Consequently, an
unbounded payoff or a finite error decomposition alone does not establish
novelty. Quantum signal-processing pricing and estimation are further
published opponents \citep{stamatopoulos,rallfuller2023}. I compare a bounded Chebyshev sampling adaptation on all three selected
first-correction preparations in each of four cases, with parity-correct
circuits, confidence spending and a charged fallback. It does not reproduce every heuristic of published ChebAE,
replace the financial payoff arithmetic, or defeat an unimplemented
upstream QSP evaluator.

The common-information Gaussian bound follows established convolution and
elementary calculus. Common-input sensitivity coupling is longstanding
\citep{glassermanyao1992}, while \citet{chakrabarti} already separates
intrinsic and inherited quantum-financial arithmetic errors and propagates
derivative-based allowances. Stieltjes pricing and left-limit calculus are
established \citep{marinelli2024}; its call-spanning measure differs from
the bounded increasing payoff measure used here. I therefore claim neither
a new generic stability principle nor a new quantum sampling law.
The reflection-cost lemma specialises the query and parity constructions
of \citet{suzuki2020,rallfuller2023}. Cross-boundary circuit simplification
is established \citep{stamatopoulos2020,carreravazquez2021}.
\footnote{The source ledger identifies exact consulted sections and versions.
For Carrera Vazquez--Woerner, final APS metadata and abstract were verified,
but detailed Section~IV was read in arXiv:2005.07711v1, dated 15 May 2020.
I do not imply that the inaccessible final theorem text was read.}
The contribution candidate is the exact guarded common-endpoint allocation,
its reusable sharp component envelopes and the comparative financial
consequences after the same optimisation is supplied to every rival.

The project-specific argument combines sparse payoff-boundary curvature,
shared certified norms and positive or signed coherent corrections with a
moment-factorised Gaussian transport certificate. The latter integrates
derivative envelopes against positive error measures, keeping the small
basket-dispersion terms which a global box maximum loses. I derive each part and test
their complete composition against stronger evaluators. The general
positive-threshold norm theorem is distinguished from the sharper
payoff-specific result. The positive second-order cell identity is an exact specialisation of
\citet{cerone2002}, Theorem~3, equations~(2.13)--(2.15). Neither that
identity, its transform nor block refinement is new generic theory.
Transform-kernel representations have further published antecedents
\citep{blair2016}.\footnote{Only the publisher abstract of Blair et al. was accessible; its full proof-level comparison remains incomplete.}
The narrower contribution candidate is the residual-specific derivative
envelope and its finite-moment closure, tested with the same information
in both transport orders. A close unpublished supplement
\citep{herman2026} bounds GBM discretisation on uniform physical-coordinate
left-endpoint grids. I adapt its first-derivative mechanism to the actual
rounded equal-CDF grid; its asymptotic theorem is not a numerical guarantee
for these different nodes. The completed common-grid comparison identifies
smaller second-order allowances, but almost no value from block refinement.
The numerical comparison concerns a frozen compiler,
a declared finite precision pool and synthetic models. It can
identify conditional implementation gains and adverse cases; it cannot
prove a globally optimal quantum algorithm, theoretical priority, a device
advantage or empirical exotic-price accuracy. Detailed proofs, consulted
source versions and earlier rejected investigations are retained outside
the main argument.

\section{Financial model and the frozen contracts}
\label{sec:model}

Let $\tau>0$, $r_f\in\R$ and dividend yields $q_i\in\R$ be constant.
Under a physical probability, suppose the ex-dividend stock system is
\begin{equation}
 dS_i/S_i=b_i\,dt+\Sigma_i\,dW,
\end{equation}
where the filtration is generated by a $d_B$-dimensional Brownian motion
and the constant matrix $\Sigma\in\R^{n\times d_B}$ has full row rank.
Choose $\theta$ satisfying
$\Sigma\theta=b-r_f\mathbf1+q$.  The exponential density $L_t=\exp(-\theta\cdot W_t-\|\theta\|^2t/2)$
defines an equivalent measure $Q$ on the finite horizon. Under $Q$, asset
drifts are $r_f-q_i$ and discounted cumulative gains are true martingales.
Admissible self-financing discounted wealth bounded below is a
supermartingale, precluding arbitrage. Appendix~\ref{app:measure-change}
gives the Gaussian density shift, integrability and localisation proof.
\footnote{Full row rank is sufficient for the stated measure construction.
If Brownian risk is unspanned, I use the stipulated pricing measure $Q$;
the quantum estimator does not select another measure.}

At maturity, independent standard normals $Y\in\R^p$ and $Z\in\R$ give
\begin{equation}
 S_i=K\exp(\mu_i+R_i\cdot Y+\lambda Z),\qquad
 \mu_i=(r_f-q_i)\tau-\tfrac12(\|R_i\|^2+\lambda^2),
 \label{eq:stocks}
\end{equation}
with $K>0$ and $\lambda>0$.  The loadings include the horizon scale.
It\^o's formula for $\log S_i$ proves this expression and gives
$\E_QS_i=Ke^{(r_f-q_i)\tau}$.  For positive weights summing to one, define
\begin{equation}
 \begin{aligned}
 A&=\sum_iw_iS_i, &G&=\prod_iS_i^{w_i},\\
 F_H(x)&=\min\{(x-K)_+,H\}, &F_\infty(x)&=(x-K)_+,\\
 V&=\E_QF(A), &\Pi&=e^{-r_f\tau}V.
 \end{aligned}
 \label{eq:financial-target}
\end{equation}
Here $F$ is the specified capped or ordinary call. Its time-$t$ price is
$e^{-r_f(\tau-t)}\E_Q[F(A)\mid\mathcal F_t]$.
For either payoff $0\leq F(A)\leq A$ and
$\E A^2\leq\sum_iw_i\E S_i^2<\infty$ by convexity.
The discounted conditional price is therefore a true square-integrable
martingale. Adding this claim at that price preserves the same no-arbitrage
argument. Boundedness of the raw payoff is unnecessary.  \footnote{The numerical certificate concerns this single model price.
Separately rounded quotes across strikes need not form an arbitrage-free
quoted surface.}

Write
\begin{equation}
 \mu_G=\sum_iw_i\mu_i,\quad R_G=\sum_iw_iR_i,\quad
 d_i=\mu_i-\mu_G,\quad s_i=R_i-R_G,
 \quad m(y)=\mu_G+R_G\cdot y.
 \label{eq:geometry}
\end{equation}
Then $\sum_iw_id_i=0$ and $\sum_iw_is_i=0$ exactly.  With
\begin{equation}
 T(y)=\sum_iw_ie^{d_i+s_i\cdot y},\qquad U(y)=T(y)-1,
 \label{eq:ratio}
\end{equation}
the identity $A=GT(Y)$ holds.  Jensen's inequality gives $T\geq1$;
hence the arithmetic basket exceeds its geometric control pointwise.
This deterministic ordering justifies the positive first correction.
It does not imply that the ratio and the geometric level are independent.

Every positive-definite maturity log-covariance matrix $C$ admits the
common loading
\begin{equation}
\lambda^2=(\mathbf1^{\mathsf T}C^{-1}\mathbf1)^{-1},\qquad
RR^{\mathsf T}=C-\lambda^2\mathbf1\mathbf1^{\mathsf T}.
\end{equation}
The residual matrix is positive semidefinite of rank $n-1$; its factor and
an independent common normal reproduce the original covariance and
drifts. Appendix~\ref{app:common-factor} proves these claims. The
decomposition establishes model coverage, while numerical cost remains
a separate comparison.

The numerical family has $K=100$, $\tau=1$, $r_f=0.03$, $q_i=0.02$,
$n\in\{2,4,8\}$, $\lambda\in\{0.05,0.2,0.4\}$ and
$H\in\{1,10,25\}$, giving 27 capped contracts from nine covariance models.
I add an ordinary call for each model, giving 36 financial cases.
The two-asset anchor uses $w=(1/2,1/2)$ and $R=(-0.1,0.2)^{\mathsf T}$.
For four and eight assets, the residual loading matrix is a scaled Helmert
contrast matrix, with weights
$(0.4,0.3,0.2,0.1)$ and
$(0.3,0.2,0.15,0.1,0.08,0.07,0.06,0.04)$ respectively.
Its scale gives weighted contrast dispersion $0.0225$ in the ideal matrix,
matching the two-asset anchor.  The actual 70-digit decimal loadings are
frozen as exact rationals, their drifts are computed from those rationals,
and exact rational $LDL^{\mathsf T}$ factorisation checks positive
covariance pivots and full volatility row rank.  The ideal-to-frozen
dispersion discrepancy is below $1.8\times10^{-71}$.
Matching this one dispersion does not match the full payoff distribution
or isolate a causal effect of dimension.

These are reproducible synthetic experiments. All figures and tables use
synthetic model inputs or generated verification results; no LSEG
observations or inaccessible option sample enter the certificates.
\footnote{Any future LSEG calibration is restricted to my authorised
individual-use entitlement. Extraction and analysis code, together with
only the permitted limited academic, non-commercial summaries, may be
released. Underlying observations and systematic derived datasets remain
excluded from GitHub and other third-party redistribution under the
supplied institutional guidance.}
Historical returns alone would not validate an observed exotic price.
The present claims depend on exact model inputs and retained arithmetic
certificates, rather than an unreported private empirical sample.

\section{Conditional smoothing and the joint derivative certificate}
\label{sec:norm}

The capped payoff has the positive layer-cake representation
$F_H(x)=\int_K^{K+H}\ind\{x>k\}\,dk$.
Tonelli's theorem therefore gives the conditional kernel
\begin{equation}
 P_{\lambda,H}(m)=\E_ZF_H(Ke^{m+\lambda Z})
 =\int_0^{\log(1+H/K)}Ke^x\Phi((m-x)/\lambda)\,dx.
 \label{eq:kernel}
\end{equation}
The undiscounted price is
$V=\E_Y P_{\lambda,H}(m(Y)+\log T(Y))$.
Conditioning has removed the nonsmooth payoff threshold from the financial
integrand.  It also supplies the exact geometric control
$c_0=\E P_{\lambda,H}(m(Y))$.
The positive first correction
$R_1=P_{\lambda,H}(m+\log T)-P_{\lambda,H}(m)$ is in $[0,H]$.
Gaussian smoothing and control identities are useful classical reductions
before any quantum encoding is considered \citep{curran,bayer}.

The derivative result extends beyond a single capped call.  Let $\nu$ be a
finite nonnegative Borel measure on real log thresholds, of mass $H$,
and define
\begin{equation}
 P_\nu(m)=\int\Phi((m-x)/\lambda)\nu(dx),\qquad
 g_m(t)=P_\nu(m+\log t),\quad t>0.
 \label{eq:measure-kernel}
\end{equation}
For the capped call,
$\nu(dx)=Ke^x\ind_{[0,\log(1+H/K)]}(x)dx$.
A cash digital is a point mass; finite positive mixtures of thresholds are
also covered.  Negative threshold weights, unbounded payoff mass and
negative basket weights are not hypotheses of this theorem.

Let $s(r,\ell)$ be signed Stirling numbers of the first kind, with
$s(0,0)=1$ and
$s(r+1,\ell)=s(r,\ell-1)-rs(r,\ell)$.
The Hermite polynomials satisfy
$\operatorname{He}_0=1$, $\operatorname{He}_1=z$ and
$\operatorname{He}_{j+1}=z\operatorname{He}_j-j\operatorname{He}_{j-1}$.
Define
\begin{equation}
 A_{r,\lambda}(z)=\sum_{\ell=1}^r s(r,\ell)\lambda^{-\ell}
 (-1)^{\ell-1}\operatorname{He}_{\ell-1}(z),\qquad
 D_r(\lambda)=\sup_z|A_{r,\lambda}(z)\phi(z)|.
 \label{eq:joint-norm}
\end{equation}
Every polynomial times the Gaussian density is bounded and tends to zero
at both infinities.  In particular $D_r$ is finite and strictly positive.
The separate Hermite norms
$\mathcal H_j=\sup_z|\operatorname{He}_j(z)\phi(z)|$ give the valid
triangle bound
$D_r\leq\sum_{\ell=1}^r|s(r,\ell)|\lambda^{-\ell}\mathcal H_{\ell-1}$.
Computing the combined norm preserves cancellation between derivative
terms which this triangle inequality loses.

\begin{theorem}[Sharp positive-threshold bound]
\label{thm:norm}
For $r\geq1$, $\lambda>0$ and the finite positive measure in
\eqref{eq:measure-kernel},
\begin{equation}
 g_m^{(r)}(t)=t^{-r}\int
 A_{r,\lambda}((m+\log t-x)/\lambda)
 \phi((m+\log t-x)/\lambda)\nu(dx),\qquad
 |g_m^{(r)}(t)|\leq HD_r(\lambda)t^{-r}.
 \label{eq:norm-derivative}
\end{equation}
The constant $HD_r$ is the smallest uniform constant over measures of mass
$H$, real $m$ and positive $t$.
\end{theorem}
The falling log-price derivative operator produces the polynomial in \eqref{eq:joint-norm}; positivity bounds its integral. A point-mass threshold attains the resulting operator norm. Appendix~\ref{app:proof-norm} gives the full differentiation and sharpness proof.

This sharpness concerns the stated measure class; a particular capped-call
measure or basket remainder need not attain the norm.  The proof also
identifies exactly which object must be computed.  For rational $\lambda$,
$Q_{r,\lambda}=A_{r,\lambda}'-zA_{r,\lambda}$ is a rational polynomial of
degree $r$.  Every positive absolute maximum of $A_{r,\lambda}\phi$ lies
at a real root of $Q_{r,\lambda}$.  Exact square-free reduction, a rational
root-counting sequence and finite rational isolation find every distinct
root; Gaussian evaluations over the isolating intervals give an outward
upper certificate $\overline D_r\geq D_r$.
Appendix~\ref{app:roots} supplies the termination and enclosure argument.
All finite calculations below use $\overline D_r$, rather than an
uncertified maximisation on a numerical mesh.  Their arithmetic is
implemented with Arb midpoint-radius enclosures \citep{johansson2017}.

The separate Hermite and Fourier triangle norms remain competent comparators. They distinguish cancellation gains from tighter individual derivative bounds; Appendix~\ref{app:fourier-comparators} gives their recurrences.

\section{Exact Gaussian coefficients and finite classical screening}
\label{sec:coefficients}

Put $\sigma^2=\lambda^2+\|R_G\|^2$.
For a multi-index $\gamma\in\mathbb N_0^n$, define
$|\gamma|=\sum_i\gamma_i$, $\gamma!=\prod_i\gamma_i!$,
$w^\gamma=\prod_iw_i^{\gamma_i}$,
$d_\gamma=\sum_i\gamma_id_i$ and
$v_\gamma=\sum_i\gamma_is_i$.
The multinomial theorem and completion of the Gaussian square imply
\begin{align}
 T(y)^j&=\sum_{|\gamma|=j}\frac{j!}{\gamma!}w^\gamma
                   e^{d_\gamma+v_\gamma\cdot y},\label{eq:T-multinomial}\\
 \E[e^{v\cdot Y}P_{\lambda,H}^{(\ell)}(\mu_G+R_G\cdot Y)]
 &=e^{\|v\|^2/2}P_{\sigma,H}^{(\ell)}(\mu_G+R_G\cdot v).
 \label{eq:gaussian-tilt}
\end{align}
The square-completion and dominated-differentiation argument is given in Appendix~\ref{app:gaussian-tilt-detail}.

For ordinary calls the same tilt formula follows from the Gaussian
exponential moment and dominated differentiation on compact argument
intervals. In the second-correction study I need only
\begin{equation}
 c_1=\sum_iw_i e^{d_i+\|s_i\|^2/2}
 P_\sigma'(\mu_G+R_G\cdot s_i)-P_\sigma'(\mu_G),
\end{equation}
where $P_\sigma$ is the chosen capped or vanilla kernel.
The finite positive-threshold theorem below remains a capped-payoff result;
the vanilla remainder uses the boundary theorem in Section~\ref{sec:oracles}.
The shift $R_G\cdot v$ retains the covariance between the geometric level
and the ratio.  Removing it generally changes the coefficient.

Let
\begin{align}
 J_{j,\ell}&=\sum_{|\gamma|=j}\frac{j!}{\gamma!}w^\gamma
 e^{d_\gamma+\|v_\gamma\|^2/2}
 P_{\sigma,H}^{(\ell)}(\mu_G+R_G\cdot v_\gamma),\label{eq:J}\\
 \mathcal M_r&=\sum_{j=0}^r(-1)^{r-j}\binom rj
 \sum_{|\gamma|=j}\frac{j!}{\gamma!}w^\gamma
 e^{d_\gamma+\|v_\gamma\|^2/2}.
 \label{eq:continuous-moments}
\end{align}
Then $\mathcal M_r=\E U^r$ and all terms are finite.
Taylor's theorem about $t=1$ defines
\begin{equation}
 c_0=P_{\sigma,H}(\mu_G),\qquad
 c_a=\frac1{a!}\sum_{\ell=1}^as(a,\ell)
        \sum_{j=0}^a(-1)^{a-j}\binom ajJ_{j,\ell},\quad a\geq1.
 \label{eq:coefficients}
\end{equation}
The inner alternating sum is the exact expectation of
$U^aP_{\lambda,H}^{(\ell)}(m(Y))$.
Its cancellation is kept inside outward arithmetic, not estimated from
printed decimal precision.

\begin{theorem}[Finite Gaussian price certificate]
\label{thm:finite-price}
Under \eqref{eq:stocks}--\eqref{eq:ratio}, for every finite integer $k\geq0$,
\begin{equation}
 \left|V-\sum_{a=0}^kc_a\right|
 \leq\frac{H\overline D_{k+1}(\lambda)\mathcal M_{k+1}}{(k+1)!}.
 \label{eq:finite-certificate}
\end{equation}
Each coefficient and moment is evaluated by the finite sums above.
For rational inputs and positive common scale, validated enclosures of all
endpoints exist at any prescribed positive arithmetic accuracy.
\end{theorem}
The integral Taylor remainder and Theorem~\ref{thm:norm} give the displayed radius. Gaussian exponential moments make every finite-order coefficient and remainder integrable. Appendix~\ref{app:proof-finite-price} supplies the full argument.

For ordinary calls and finite $k\geq1$, write $C_r=\prod_{j=0}^{r-1}(\partial_m-j)P_{\lambda,\infty}$ and take $M_r\geq\|C_r\|_\infty$ for $r\geq2$. The same Gaussian tilts evaluate
$c_a=\E[U^aC_a(m)]/a!$ and yield
\begin{equation}
 \left|V-\sum_{a=0}^kc_a\right|
 \leq M_{k+1}\E U^{k+1}/(k+1)!.
\end{equation}
Indeed, $g_m^{(k+1)}(t)=t^{-k-1}C_{k+1}(m+\log t)$ is bounded by
$M_{k+1}$ for $t\geq1$. The integral Taylor remainder and finite Gaussian
moments prove the inequality exactly as above. The unbounded $P,P'$ terms
remain integrable under each Gaussian tilt; the higher falling derivatives
have bounded Gaussian boundary kernels. At $k=1$ positivity of the ordinary
call curvature gives the one-sided interval in
Corollary~\ref{cor:convex-screen}. This is a finite-order classical
competitor, with no infinite-series convergence claim.

If $[L_k,U_k]$ encloses the coefficient sum and $\overline R_k$ bounds the
right side, I use
$[L_k-\overline R_k,U_k+\overline R_k]\cap[0,H]$.
Intersections across finitely many retained degrees can only improve this
interval.  I stop when its midpoint radius meets the monetary target,
and otherwise use a separately certified method or declare failure.
The argument is finite-order: finite lognormal moments do not establish
convergence of the infinite Taylor expansion.  The earlier Fourier
triangle remainder eventually grows for any nonzero centred loading,
so taking ever more terms is not a valid convergence policy.

The kernel evaluations are explicit.  With $L_H=\log(1+H/K)$,
$x_0=m/\sigma$, $x_1=(m-L_H)/\sigma$ and $z_i=x_i+\sigma$,
\begin{equation}
 P_{\sigma,H}(m)=K\{e^{m+\sigma^2/2}[\Phi(z_0)-\Phi(z_1)]
                   -\Phi(x_0)+(1+H/K)\Phi(x_1)\}.
 \label{eq:closed-kernel}
\end{equation}
Independent derivative-series and strike-integral checks verify the kernel evaluations; Appendix~\ref{app:kernel-differentiation} gives the cancellation identities and differentiation details.

The certificate extends to unequal positive spots by replacing $d_i$ with
$\log(S_{i,0}/G_0)+\mu_i-\mu_G$, where
$G_0=\prod_iS_{i,0}^{w_i}$, and shifting every scalar-kernel argument by
$\log(G_0/K)$.  Weighted centring and the covariance shift survive.
Positive asset-monitor observations can likewise be treated as basket
components when their joint Gaussian log prices admit an equal independent
common loading $\lambda>0$.  The component count then includes monitoring
times; the multinomial cost can grow substantially.  These are analytical
extensions, not additional empirical tests.

\section{Payoff curvature and coherent correction evaluators}
\label{sec:oracles}

Let $P_F(m)=\E F(S_{\rm ref}e^{m+\lambda Z})$, $D=\partial_m$ and
$C_r=D(D-1)\cdots(D-r+1)P_F$. I write $m=m(Y)$ and $U=U(Y)$.
Taylor's integral identity gives the exact correction
\begin{equation}
 \begin{aligned}
 R_2&=P_F(m+\log(1+U))-P_F(m)-UP_F'(m)\\
 &=\frac{U^2}{2}\E_{Z_B}[(1+Z_BU)^{-2}C_2(m+\log(1+Z_BU))].
 \end{aligned}
 \label{eq:beta-identity}
\end{equation}
where $Z_B\sim\operatorname{Beta}(1,2)$. The zero extension at $U=0$
requires no division. Analytic evaluation of the payoff boundary removes
the strike auxiliary variable from the earlier integral construction.

\begin{theorem}[Sparse payoff-boundary correction]
\label{thm:boundary}
Suppose $F(s)=a+bs+\sum_{j=1}^J\alpha_j(s-K_j)_+$ with finite $J$ and
positive $K_j,S_{\rm ref},\lambda$. Then
\begin{equation}
 C_2(m)=\sum_{j=1}^J\frac{\alpha_jK_j}{\lambda}
 \phi\!\left(\frac{m-\log(K_j/S_{\rm ref})}{\lambda}\right).
 \label{eq:boundary-kernel}
\end{equation}
If $\alpha_j\geq0$ for all $j$, then $R_2\geq0$ for every $U\geq0$.
For all finite signed kink weights,
$|R_2|\leq M_2U^2/2$ whenever $M_2\geq\|C_2\|_\infty$.
\end{theorem}
Each payoff kink contributes a Gaussian boundary density to $P_F^{\prime\prime}-P_F^\prime$; the affine payoff terms vanish under $D(D-1)$. The nonnegative Taylor weight gives the sign under convexity and the absolute norm bound for signed kink weights. Appendix~\ref{app:proof-boundary} proves every step.

Adding $a+bs$ to the payoff leaves $R_2$ exactly unchanged: its
conditional $g_m(t)=a+bS_{\rm ref}e^{m+\lambda^2/2}t$ is affine in $t$
and has zero second remainder. Calls and puts at the same strike therefore
share the ideal curvature correction; the deterministic coefficients
restore their established discounted difference
$e^{-r_f\tau}(\E A-K)$. This is an applicability consequence of the
derivation, rather than a new put--call parity claim or an additional
numerically tested payoff family.

\begin{corollary}[Continuous convex screening]
\label{cor:convex-screen}
For a convex finite-kink payoff under \eqref{eq:stocks}, suppose
$M_2\geq\|C_2\|_\infty$. Then, in every retained asset dimension,
\begin{equation}
 c_0+c_1\leq V\leq c_0+c_1+M_2\E U^2/2.
\end{equation}
\end{corollary}
Taking expectations of the positive bounded residual gives the interval, with Gaussian tilts evaluating its coefficients. Appendix~\ref{app:proof-convex-screen} gives the complete proof.
This analytical screen precedes financial discretisation and quantum
estimation. It does not require an exponentially large product grid or
assume that a small correction deserves a quantum computation.

For the cap, set $c=1+H/K$, $\ell=\log c$ and $z=m/\lambda$; then
\begin{equation}
 C_2(m)=\frac K\lambda\{\phi(z)-c\phi(z-\ell/\lambda)\}.
 \label{eq:cap-curvature}
\end{equation}
For the vanilla call $C_2=K\phi(m/\lambda)/\lambda\geq0$ and its exact
norm is $K/(\lambda\sqrt{2\pi})$. I use a rational outward upper endpoint
$M_2$ in every graph. For the cap, the two stationary points in
Appendix~\ref{app:payoff-norms} certify its exact global norm. I also
certify the unique global peak of $P'$ as $M_1$.
At $\lambda=0.05,H=25$, the resulting upper constants are approximately
39.23 and 54.99 per cent of $H\overline D_2$ and $H\overline D_1$,
respectively. The first-derivative comparator also has the valid bound
$\min\{K+H,H\overline D_1\}$; against that stronger bound, its reduction
at this case is approximately 12.25 per cent. These constants are supplied
to every current coherent opponent and
the new classical evaluators; they are not reserved for the proposed oracle.

On the realised financial support, $0\leq U\leq U_*$.
For $U_*>0$, define
\begin{equation}
 \begin{aligned}
 B_2&=M_2U_*^2/2, &t&=1+z_BU,\\
 a&=(U/U_*)^2t^{-2}, &h&=(1+C_2(m+\log t)/M_2)/2.
 \end{aligned}
 \label{eq:factorised-kernel}
\end{equation}
For capped calls, the boundary evaluator prepares
\begin{equation}
 p^{\rm boundary}=(1-a)/2+ah,\quad X=2p^{\rm boundary}-1,
 \quad B_2X=U^2t^{-2}C_2(m+\log t)/2.
 \label{eq:integral-probability}
\end{equation}
All factors lie in $[0,1]$ and $|X|\leq U^2/U_*^2$.
For vanilla calls the native positive encoding is
$p^{\rm boundary}=a\gamma(m+\log t)$ with
$\gamma=C_2/M_2\in[0,1]$, and its monetary decoder is $B_2p$, rather
than $B_2(2p-1)$. The raw payoff is unbounded, but this encoded correction
is bounded on the certified finite financial support. Gaussian tails
remain a separate financial obligation.

A strong deterministic evaluator uses the two-node Gauss-Jacobi rule for
$2(1-z)dz$:
\begin{equation}
 z_\mp=(4\mp\sqrt6)/10,\quad
 w_\mp=\tfrac12\pm1/(3\sqrt6),\quad
 R_2^{\rm GJ}=\frac{U^2}{2}\sum_\pm
 w_\pm(1+z_\pm U)^{-2}C_2(m+\log(1+z_\pm U)).
 \label{eq:stable-kernel}
\end{equation}
It preserves the same small-factor representation without a sampled
auxiliary or raw price subtraction. I implement its exact dyadic constants,
with complementary positive weights summing to one, and charge the
algebraic node and weight discrepancies. A fourth-order Hermite remainder
proves its monetary quadrature bound, not a heuristic comparison of two
quadrature values.

The raw direct opponent uses the same $c_0+c_1$, $M_2$ and $B_2$:
\begin{equation}
 p^{\rm raw}=\tfrac12+R_2/(2B_2)
 \label{eq:direct-probability}
\end{equation}
for the cap and $p^{\rm raw}=R_2/B_2$ for vanilla calls. Its arithmetic
errors are propagated before cancellation. The positive first-correction
opponent uses $c_0$, $R_1=P(m+\log T)-P(m)$ and
\begin{equation}
 p^{(1)}=R_1/B_1,\qquad B_1=\min\{H,M_1U_*\}
 \label{eq:first-probability}
\end{equation}
for the cap. For vanilla calls $B_1=M_{1,\rm finite}U_*$, where
$M_{1,\rm finite}=Ke^{m_++\lambda^2/2}$ bounds the conditional forward on
its certified argument interval. There is no globally bounded vanilla
first log-price derivative. Every scalar, normaliser and input-domain
certificate is method-specific, with identical information sharing.

\begin{proposition}[Auxiliary variance penalty]
\label{prop:variance}
With continuous auxiliaries and the same finite financial law, the exact
boundary score has conditional mean $R_2(Y)$ and
$\Var(B_2X)=\Var(R_2(Y))+\E\Var(B_2X\mid Y)$ for the capped encoding;
the analogous identity uses $B_2p$ for the vanilla encoding.
\end{proposition}
The conditional mean is $R_2(Y)$, so the law of total variance separates financial dispersion from auxiliary randomisation. The full square-expansion argument is in Appendix~\ref{app:proof-variance}.
This established variance decomposition prevents attributing an intrinsic
variance reduction to auxiliary randomisation. Any gain must survive its
actual evaluator and estimation costs. The stable quadrature opponent is
an implemented comparator; the finite pool is not a lower bound on every
possible cancellation-preserving pricing method.

\section{Finite preparation and complete monetary error}
\label{sec:finite}

Independent uniform labels prepare a realised truncated-normal quantile
law. With $L>0$, $N=2^{b_y}$ and $\alpha_1=\Phi(L)-\Phi(-L)$, I use
\begin{equation}
 y_j=\Phi^{-1}(\Phi(-L)+\alpha_1(j+1/2)/N),\quad
 \widehat y_j=2^{-q}\operatorname{round}(2^qy_j),\quad \alpha=\alpha_1^p.
 \label{eq:quantile-grid}
\end{equation}
Outward CDF brackets prove unique rounding, order and domain membership.
Duplicate coordinates keep their label multiplicities. Hadamards prepare
labels; coherent lookup computes coordinates and its inverse clears work.
If $[\ell_y,h_y]$ covers every ideal enclosure and actual node, then
\begin{equation}
 U_*\geq\max_{v\in\{\ell_y,h_y\}^p}\sum_iw_ie^{d_i+s_i\cdot v}-1
 \label{eq:finite-vertex}
\end{equation}
is sufficient: the sum is convex in each coordinate, so maximisation can
move successively to endpoints. Outward vertex evaluation certifies all
nine model bounds. The continuous tail box and finite normalisation support
have different roles.

The boundary method samples the exact cell-mass midpoint Beta law. The
minimum of two $b_Z$-bit uniform labels has mass $2h_Z(1-z_j)$ at
$z_j=(j+1/2)h_Z$, $h_Z=2^{-b_Z}$, which equals the continuous
$\operatorname{Beta}(1,2)$ mass of that cell. No strike register is needed.
The stable method instead uses the two exact dyadic Gauss constants.
Let $M_{r,N}=\E_NU^r$ under the actual financial label law. Appendix~\ref{app:aux}
proves, before the external financial factor,
\begin{align}
 \beta_{\rm B}&\leq h_Z^2(M_3M_{3,N}/12+M_4M_{4,N}/48),
 \label{eq:beta-error}\\
 \beta_{\rm GJ}&\leq M_6M_{6,N}/14400+
 \frac{M_3M_{3,N}}2\sum_\pm w_\pm\delta z_\pm+
 \frac{M_2M_{2,N}}2\sum_\pm\delta w_\pm.
 \label{eq:gauss-error}
\end{align}
Here $M_r\geq\|C_r\|_\infty$ is a derivative bound, distinct from the
moment $M_{r,N}$. The common norm ledger uses
$C_{r+1}=(D-r)C_r$, exact polynomial stationary-root isolation and outward
Gaussian evaluation. Actual fourth and sixth moments strengthen the
stable opponent; finite-support fallbacks remain valid but are not used
to weaken it. Raw direct and first-correction methods have zero auxiliary
error.

Every function is implemented by a one-dimensional endpoint table with
outward nodal enclosures and signed-magnitude interpolation. The graph
propagates absolute errors and implemented ranges through exponentials,
positive ratio arithmetic, logarithms, inverse squares and the specified
payoff kernels. Direct subtraction receives no assumed favourable relative
error. Domain violations, overflow and ambiguous signed ranges reject a
graph. Appendix~\ref{app:arithmetic} states the interpolation and arithmetic
contracts; the supporting source retains every node. Computing a numerical
threshold, toggling a comparator flag and applying the literal inverse
cleans private work. For integer $k\in\{0,\ldots,2^q\}$ and uniform
$j\in\{0,\ldots,2^q-1\}$, the predicate $j<k$ gives exact Born success
probability $k/2^q$, including both endpoints. Complete emitted preparations
replace any free-oracle assumption.

Appendix~\ref{app:financial-transfer} derives the price-independent
financial certificate. A local signed-cell calculation supplies the
sharper two-asset bounds. For higher dimensions, I express the cell error
as a positive measure, compute its exponential moments and use tensor
hybrids to integrate dispersion-weighted derivative envelopes. All 36
first- and second-correction targets now have complete financial transfers.
The global transport family retains its three insufficient two-asset bounds;
I use the stronger local certificate in those cases. The construction is
classical setup, with explicit operation counts and outward errors.

\begin{theorem}[Complete monetary error]
\label{thm:error}
Let $e_c$ bound the chosen coefficient midpoint error, let
$[F_-,F_+]$ enclose the external mass with midpoint and radius $e_F$, and
let the continuous-to-finite correction discrepancy be at most $\beta_Y$.
Suppose the ideal and implemented probabilities differ by at most $\eta_p$
and the implemented score has absolute value at most one. On the jointly
certified estimation event,
\begin{equation}
 |\widehat V-V|\leq e_c+\beta_Y+
 F_+(\beta_{\rm aux}+dB\eta_p+e_{\rm stat})+e_FB,
 \label{eq:total-error}
\end{equation}
where $(B,d)=(B_2,2)$ for a capped centred second correction and
$(B,d)=(B_2,1)$ for a native positive second correction, or $(B_1,1)$ for
the first correction. The auxiliary term is chosen from
\eqref{eq:beta-error} or \eqref{eq:gauss-error}, or is zero for a raw
opponent. Statistical error is monetary before the mass, at its separately
certified joint confidence. Discounting multiplies the bound by
$e^{-r_f\tau}$.
\end{theorem}
Each deterministic replacement is charged before the bounded probability decoder is applied. The decoder contributes $dB\eta_p$, and uncertain external mass contributes $e_FB$ on the joint estimation event. Appendix~\ref{app:proof-error} gives the full decomposition.

A graph is admitted by the complete monetary budget in
Theorem~\ref{thm:error}, rather than by its number of fractional bits.

\section{Signed flags and competent finite moment priors}
\label{sec:priors}

The centred encoding can have probability near one half even when its
monetary correction is small.  I therefore also implement two exact signed
parts of the represented probability:
\begin{equation}
 \widehat p_+=\max(0,2\widehat p-1),\qquad
 \widehat p_-=\max(0,1-2\widehat p),\qquad
 \widehat p_+-\widehat p_-=2\widehat p-1.
 \label{eq:signed-parts}
\end{equation}
The transformations use exact dyadic scaling and positive subtraction.
They preserve the difference on every represented basis input, including
$\widehat p\in\{0,1/2,1\}$.
They are separate success flags with separately charged preparation and
confidence costs; I do not estimate a fictitious flag with the variance of
a signed classical output.  Their deterministic recombination uses the
centred base certificate in Theorem~\ref{thm:error}, avoiding a double
charge for an exact integer identity.

The retained truncated-normal quantile law has $\omega=1$ and external
mass $\alpha=\alpha_1^p$. For the uniform-box variant on $[-L,L]^p$,
$\omega(y)=e^{-\|y\|^2/2}$ and $C_L=(2L/\sqrt{2\pi})^p$.
These weights belong to the prepared score; their monetary factors are
decoded after estimation. The uniform-box variant prepares
$p_\omega=(1+X_\omega)/2$, where $X_\omega=\omega X$, and forms the
signed parts of its represented probability in \eqref{eq:signed-parts}.
The quantile variant has $X_\omega=X$. For each capped construction,
$p_\pm\leq|X_\omega|\leq\omega U^2/U_*^2$.
The positive-part map is one-Lipschitz, so the implemented flags satisfy
\begin{equation}
 \E_N\widehat p_\pm\leq
 u_2:=\min\{1,M_{2,N}/U_*^2+2\eta_p\}.
 \label{eq:rare-prior}
\end{equation}
The expectation is under the prepared label law.  Neither its Gaussian
box mass $\alpha$ nor its uniform-box monetary factor $C_L$ belongs in
this probability prior.  They enter monetary reconstruction after
estimation.  A prior reduction improves a canonical estimator's amplitude
error bound only if the increased flag-preparation costs and two confidence
allocations are also included.

For the native positive vanilla correction the tighter implemented prior is
$u_2=\min\{1,M_{2,N}/U_*^2+\eta_p\}$, since its decoder has no factor two.
These bounds also apply to the stable rule: its positive weights sum one
and preserve the same pointwise envelope.

The finite moments are exactly computable without a financial price.
Define the one-dimensional exponential sum
\begin{equation}
 \mathcal L_N(v)=N^{-1}\sum_{j=0}^{N-1}e^{v\widehat y_j},
 \label{eq:finite-mgf}
\end{equation}
and, for uniform-box weighting, replace each summand by
$e^{vy_j-y_j^2/2}$.  Independence and the multinomial expansion give
\begin{equation}
 \E_N[\omega T^r]=
 \sum_{|\gamma|=r}\frac{r!}{\gamma!}w^\gamma e^{d_\gamma}
 \prod_{k=1}^p\mathcal L_N\!\left(\sum_i\gamma_is_{ik}\right),\qquad
 M_{r,N}=\sum_{j=0}^r(-1)^{r-j}\binom rj\E_N[\omega T^j].
 \label{eq:finite-moment}
\end{equation}
These are finite sums over actual dyadic nodes.
For the weighted law the zeroth raw moment is $\E_N\omega$, not one.
All cancellations remain enclosed outwards.
Thus a moment calculation can use products of one-dimensional sums instead
of $N^p$ financial price evaluations.  Its scalar preprocessing cost is
still retained; separability does not make it free.

The positive first-correction comparator also receives a competent prior.
Both tested payoffs are increasing and one-Lipschitz in the underlying price.
For fixed $Y=y$, $A=GT$ and the independent common factor therefore imply
\begin{equation}
 0\leq R_1(y)\leq Ke^{m(y)+\lambda^2/2}U(y).
 \label{eq:underlying-gap}
\end{equation}
The corresponding finite expectation retains the level-ratio correlation:
\begin{equation}
 \Gamma_N=Ke^{\mu_G+\lambda^2/2}
 \left\{\sum_iw_ie^{d_i}\prod_{k=1}^p\mathcal L_N(R_{G,k}+s_{ik})
       -\prod_{k=1}^p\mathcal L_N(R_{G,k})\right\}.
 \label{eq:correlated-gap}
\end{equation}
Expanding $e^{R_G\cdot Y}(T-1)$ proves this identity.
Combining it with $R_1\leq H$ and $R_1\leq M_1U$ gives
\begin{equation}
 u_1=\min\left\{1,
 \frac{\min\{H\E_N\omega,M_1M_{1,N},\overline\Gamma_N\}}{B_1}
 +\eta_{p,1}\right\},
 \label{eq:first-prior}
\end{equation}
where $\overline\Gamma_N$ is an outward upper endpoint.
For vanilla calls the mean bound is $\overline\Gamma_N$; I omit the
finite-cap bounds rather than assuming a globally bounded $P'$.
For the three one-dimensional grids, I also enclose $M_{4,N}$ and $M_{6,N}$
by direct nonnegative summation over the actual 1,024 nodes, charging its
2,048 conditional exponential evaluations per model.
Replacing $\E(e^{R_G\cdot Y}U)$ by a product of separate expectations
would generally invalidate this bound.  The prior is derived from moments
and the Lipschitz property, not a secretly supplied option price.

The first-correction comparator also receives the Gaussian coefficient
$c_1$, already supplied to every second-correction method. Define
$B_Q=M_2\E_Q U^2/2$. For a cap,
$\E_Q R_1=c_1+\E_Q R_2\in[c_{1,-}-B_Q,c_{1,+}+B_Q]$;
for an ordinary call the lower endpoint is $c_{1,-}$, since $R_2\geq0$.
If this continuous interval is $[L_Q,U_Q]$ and the first-residual transfer
is $\beta_1$, with $0<\alpha_-\leq\alpha\leq\alpha_+$, then
\begin{equation}
 L_N=\frac{\max\{0,L_Q-\beta_1\}}{\alpha_+},\qquad
 U_N=\min\{B_1,(U_Q+\beta_1)/\alpha_-,U_{N,\rm previous}\}.
 \label{eq:informed-first-mean}
\end{equation}
These endpoints enclose the exact finite mean: insert
$|\E_QR_1-\alpha\E_NR_1|\leq\beta_1$, use positivity and divide by
the adverse mass endpoints. Intersecting the earlier upper bound is valid.
The implemented probability mean therefore belongs to
\begin{equation}
 [\ell,u]=[\max\{0,L_N/B_1-\eta_p\},
           \min\{1,U_N/B_1+\eta_p,u_{1,\rm previous}\}].
\end{equation}
Its midpoint gives an analytical stop when
$\alpha_+B_1(u-\ell)/2$ plus the full deterministic allowance meets the
target. Otherwise $u$ supplies the canonical Bernoulli variance bound.
The full Gaussian second moment is another finite exponential sum; no
basket-price answer enters this information refinement.

Before any method estimates a correction, I also apply the same continuous
low-order price screen. The capped interval is
$[c_0+c_1-B_Q,c_0+c_1+B_Q]\cap[0,H]$; the ordinary-call interval is
$[c_0+c_1,c_0+c_1+B_Q]\cap[0,\infty)$, with all coefficient endpoints
outward. The global remainder theorem proves containment and its rational
midpoint radius proves the stop. A sufficient continuous interval requires
no grid or quantum circuit, so every competing method receives this same
decision. Higher-order classical pricing remains a separately reported
practical opponent. The supporting proof retains all interval and decoder
steps of the information refinement.

\section{Estimation and complete logical resource accounting}
\label{sec:estimation}

Let $\mathcal A$ be a clean preparation with success probability $p$,
including every financial and comparator register.  The Grover iterate is
formed from $\mathcal A$, its inverse and the good-state and zero-state
reflections.  Canonical phase estimation uses controlled powers of this
iterate \citep{brassard}; I retain the controlled global signs of the
reflections, since a global sign becomes a relative phase under control.
For $M=2^m$, the standard amplitude estimate obeys, on its success event,
\begin{equation}
 |\widetilde p-p|\leq
 \frac{2\pi\sqrt{p(1-p)}}M+\frac{\pi^2}{M^2}.
 \label{eq:canonical-error}
\end{equation}
Appendix~\ref{app:canonical-sketch} derives the bound from the two Grover eigenphases and the phase-grid displacement. The retained confidence calculation supplies its certified success probability.

If $p\leq u$, the variance term is bounded by $u(1-u)$ for $u\leq1/2$
and by $1/4$ otherwise.  I apply \eqref{eq:rare-prior} or
\eqref{eq:first-prior}, rather than assuming an arbitrary favourable
amplitude.  Exact finite postprocessing error and the compiled inverse
Fourier transform's implementation error are charged separately.
Fresh median repetitions amplify confidence by the exact binomial tail
for their certified single-shot success probability.  The allocation is
fixed before observing results.  Two signed-part estimates share a union
budget; at most 36 prespecified prices share an overall failure allowance of
one per cent.  The retained protocol assigns $1/3600$ per selected price.
\footnote{The four-case stressed comparison is reported separately from
the original 36-contract family, using the same per-price allowance.}
Deterministic selection among certified implementations requires no
statistical selection penalty because no measurement outcomes are used to
choose their resource plans.

The inverse Fourier transform uses explicitly certified Clifford and T sequences \citep{rossselinger2016}. Its retained and omitted phase errors are charged in every complete plan. Appendix~\ref{app:fourier-diagnostics} gives the synthesis details and largest-width diagnostics.

A complete count includes the initial preparation, every controlled
financial preparation and inverse, good and zero reflections, inverse
Fourier transform, confidence repetitions, observed qubits and reset of
reused shot registers.  Each method receives both coherent lookup architectures and can select
its lowest complete cap before any estimator shots.  The study retains both exact seven-T and
coherent eight-T construction profiles and their emitted chronological
tapes.  QROM mechanisms have established predecessors \citep{babbush};
using them correctly is an implementation obligation, not a novelty claim.
An exact lookup compute/erase pair costs 14 T gates, while the coherent
AND pair costs eight; ordinary arithmetic uses literal seven-T Toffolis
in both modes. Financial table generation, norm generation, finite moment sums and
financial-grid certification are reported as classical setup.
A logical T count is not directly equated to CPU time or fault-tolerant
wall-clock time.

For a centred second-correction flag with probability tolerance $e_p$,
the monetary statistical allowance is $2B_2e_p$ before the factor.
For two signed flags with tolerances $e_+,e_-$ it is
$B_2(e_++e_-)$, at their joint confidence.
The positive first correction uses $B_1e_p$, and the native positive
vanilla second correction uses $B_2e_p$.
I select only configurations satisfying the complete monetary constraint,
with at most half the requested radius assigned to deterministic error.
The final study fixes $q=28$, $b_y=10$ and the remaining Beta width
at eight, then tunes each scalar table independently within six to twelve
address bits. The unchanged all-twelve program remains in every method
pool, alongside three greedy searches with deterministic allowances
$0.001$, $0.0025$ and $0.005$. The search ranks certified cell savings
per marginal monetary error, with a fixed node-name tie rule; the final
choice minimises the actual complete T cap. Every proposed map, guard
result, chosen step and cost is retained. I also admit an exact integer
constant-propagation pass and the unchanged emitter equally for all methods.
The pass preserves represented outputs, widths and all input registers.
It produces no rewrite in the 204 higher-dimensional production graph
structures; its small-fixture simplifications are not extrapolated into
financial resource savings. This is a reproducible finite
search, rather than a global precision optimum.  Equal financial accuracy, information, confidence,
architecture and admissibility are the matched requirements.

\section{Retained findings and comparison boundary}
\label{sec:results}

The earlier correction-family study retains 502 candidate programs and
141,676 recorded allocation decisions; later registered pools are reported
separately below. The higher-dimensional pool adds 329 candidates,
117,871 genuine precision proposals and 33 infeasible starting allocations.
The combined record preserves every earlier implementation and strengthens
113 first-correction information intervals. Both coherent architectures
receive the same controls, financial accuracy and selection policy.
The family confidence allocation remains $1/3600$ per selected price,
with a union allowance of one per cent across all 36 cases.

Table~\ref{tab:transfer} identifies the resolved financial gap. The nominal
product grid contains $2^{70}$ labels in seven dimensions, but the transfer
construction evaluates finite exponential-polynomial moments, not that
Cartesian set. The derivative bounds retain $U$, $T_j$ and $T_{jj}$
moments, actual rounded nodes, Gaussian tails and all geometric-loading
terms. The loadings $R_G$ are nonzero in both higher-dimensional families;
weighted centring of $s_i$ does not justify discarding them.

\begin{table}[ht]
\centering\small
\begin{tabular}{rrrrrr}
\toprule
Assets&Outer dimension&Nominal labels&Max. $\beta_{Y,2}$&Max. $\beta_{Y,1}$&Cases\\
\midrule
2&1&$2^{10}$&$0.000937$&$0.001168$&12\\
4&3&$2^{30}$&$0.002284$&$0.004914$&12\\
8&7&$2^{70}$&$0.001107$&$0.003572$&12\\
\bottomrule
\end{tabular}
\caption{Complete continuous-to-finite financial allowances, rounded
upwards, under unchanged $L=6$, $b_y=10$, $q=28$. The target allowance is
$0.005$. Two-asset rows use the sharper local cell certificate;
higher-dimensional rows use positive transport moments. Arithmetic,
auxiliary and statistical errors are additional terms, not included in
these financial allowances. The table makes the distinction between
nominal grid size and executed certificate work explicit.}
\label{tab:transfer}
\end{table}

Shared continuous second-order intervals resolve 18 contracts before
quantum estimation. In the remaining 18 active comparisons, the best
second correction beats the informed positive first correction in ten
cases, with reductions of approximately 13.9 to 92.6 per cent.
Table~\ref{tab:n2} displays the complete-cap ratios for the entire family;
it reveals adverse cases as well as improvements. Matching one dispersion
quantity across models does not isolate a causal dimension effect.

\begin{table}[ht]
\centering\small
\begin{tabular}{rrrrrr}
\toprule
Assets&$\lambda$&$H=1$&$H=10$&$H=25$&Ordinary call\\
\midrule
2&0.05&$3.071_{\rm D}$&$3.022_{\rm D}$&$1.700_{\rm D}$&$0.349_{\rm D}$\\
2&0.2&A&$0.861_{\rm D}$&$0.689_{\rm D}$&$0.088_{\rm D}$\\
2&0.4&A&A&$0.355_{\rm D}$&A\\
4&0.05&$2.595_{\rm D}$&$1.323_{\rm D}$&$1.310_{\rm D}$&$0.261_{\rm D}$\\
4&0.2&A&A&$0.366_{\rm D}$&$0.074_{\rm D}$\\
4&0.4&A&A&A&A\\
8&0.05&A&$2.238_{\rm D}$&$1.205_{\rm D}$&$0.240_{\rm B}$\\
8&0.2&A&A&$0.584_{\rm D}$&A\\
8&0.4&A&A&A&A\\
\bottomrule
\end{tabular}
\caption{Lowest second-correction complete T cap divided by the informed
first-correction cap, rounded to three decimals. D denotes raw direct
$R_2$, B the boundary oracle and J stable Gauss-Jacobi; J is never the
lowest active second correction here. A denotes the shared analytical
price interval, with zero quantum calls for every method. A ratio below
one is a conditional logical-resource improvement, with the same monetary
radius $0.01$, confidence, architectures and precision-pool policy.
Exact caps, selected maps and full budgets remain in the archive.}
\label{tab:n2}
\end{table}

Raw direct evaluation is the cheapest second-correction implementation
in 17 active cases. For the eight-asset ordinary call at $\lambda=0.05$,
the boundary oracle has complete T cap 421,318,601,825 against
435,151,104,775 for raw evaluation, a reduction of approximately
3.18 per cent. The same case improves on the informed first correction
by about 76.0 per cent. This is a positive implementation result for
the strike-integrated, nonnegative curvature preparation in the retained
pool. Stable quadrature loses to the best second-correction implementation
in every active case; its exact finite sixth moments and fixed-parameter
error allowances are nevertheless supplied in full. The exact integer
simplifier produces no production improvement.

The independent checks reconstruct scalar contracts, integer profiles,
allocation decisions and full estimator plans, with separate small-fixture
replay of dirty numerical and flag workspaces. They verify actual node
rounding, retained information, clean inverse and controlled-phase
conventions. Large production profiles are checked by composition, rather
than exhaustive state simulation. Complete counts and setup ledgers are
retained in the supporting index and result archives.

The financial certificate uses factorised finite MGFs with explicit
exponential-series remainders. Scalar recurrence, coefficient formation,
normal quantiles, tensor-factor products and review work remain real
construction costs. The previous higher-dimensional precision search and
its audit took about 22.71 and 8.87 minutes in retained diagnostics.
These are research setup times, not single-price execution times or
conversions of T gates to CPU time.

I separately compare two discretisation certificates for the same $R_2$,
with the same actual nodes, boundary norms, dispersion envelopes and full
Gaussian tail. The first uses the positive absolute weighted
fundamental-theorem-of-calculus (FTC) measure; the second uses the
published centred Peano identity. This is a comparison of transport order,
not of the pricing corrections $R_1$ and $R_2$. Neither opponent uses a
whole-box dispersion maximum or Cartesian price enumeration.
I refine only the transported coordinate into $B=1,8,32,128,1024$
contiguous blocks, keeping other coordinates factorised. Both orders use
512-bit outward arithmetic and the same finite-moment closure.

The $0.005$ allocation concerns undiscounted financial transfer, not a
complete price or confidence guarantee. At every level, the second-order
certificate meets that allocation
in 33 of 36 cases and FTC in 32. At $B=1024$ the second-order allowances
are approximately 25.7 to 72.3 per cent smaller than FTC. Refinement from
$B=1$ tightens those second-order allowances by only 0.0702 to 0.3480
per cent and rescues none of the three two-asset failures. The previously
retained sharper local certificate still handles those cases. Thus this
experiment supports inexpensive aggregate residual certificates; it does
not identify a new refinement principle or an additional accepted price.

\begin{table}[ht]
\centering\small
\begin{tabular}{rrrrrrr}
\toprule
Assets&$p$&FTC passes&Peano passes&Refinement (\%)&$B=1$ (s)&$B=1024$ (s)\\
\midrule
2&1&9/12&9/12&0.0702--0.0839&0.019&16.95\\
4&3&11/12&12/12&0.2162--0.2573&0.095&80.13\\
8&7&12/12&12/12&0.3250--0.3481&0.470&206.67\\
\bottomrule
\end{tabular}
\caption{Value and construction cost of transported-coordinate refinement.
Pass counts concern undiscounted transfer only and are unchanged at all
five levels. Refinement is the range of
second-order allowance reductions from $B=1$ to $1024$, including the same
tail. Each time covers geometry and both certificate orders for all twelve
payoffs in that asset family, excluding common prefix setup (1.276 seconds)
and serialisation. The first level forms fresh grouped coefficients;
later levels reuse them. These joint incremental diagnostics are neither
order-specific timings nor cold-run speed ratios. The table tests whether
a tighter certificate justifies additional construction work.}
\label{tab:certificate-refinement}
\end{table}

Independent checks reproduce the one-block aggregate bounds, evaluate
38,313 production block transforms at 768-bit precision and check 198
positive interval quadratures. A separate 1,024-bit small-grid audit checks
1,584 hybrid polynomial moments. The retained finite-measure proof covers
negative and zero transform rates, raw masses and monotone Cauchy
refinement. These checks support the computed allowances; they do not
establish a uniform ordering of transport orders for other residuals.

The continuous classical test now uses only the same $M_2,M_3,M_4$
boundary constants and exponential geometry available to transport,
plus additional Gaussian moment and covariance-shifted kernel evaluations.
Finite Taylor pricing through degree three has remainder radius
$M_4\E U^4/24$. I retain the earlier successful lower orders rather than
assuming each higher order improves the bound. Table~\ref{tab:classical-information}
shows that this bounded calculation certifies 33 complete cent prices,
including all 24 higher-dimensional contracts, before preparing a grid.
Those are exactly the cases which pass the aggregate $R_2$ transfer
allocation. This is a cohort observation, not a theorem relating the
$0.005$ transfer allocation to the complete $0.01$ price radius.
Hybrid moments are not themselves continuous price coefficients; the
additional Gaussian queries and kernel evaluations are charged explicitly.

\begin{table}[ht]
\centering\small
\begin{tabular}{rrrrr}
\toprule
Assets&Degree $\leq1$&Degree $\leq2$&Degree $\leq3$&Geometric bounds\\
\midrule
2&4/12&8/12&9/12&0/12\\
4&6/12&9/12&12/12&0/12\\
8&8/12&10/12&12/12&0/12\\
\bottomrule
\end{tabular}
\caption{Complete classical cent certificates at unchanged model inputs.
Degree columns retain the union of successful orders up to the cap;
all saved outward discounted endpoints also pass the cent target.
The last column combines the published geometric exercise-set lower bound
with the arithmetic--geometric upper bound, and uses a derived separate-strike
call-spread enclosure for caps. It is not the strongest possible capped
specialisation of that method. The table identifies how much pricing work
remains when the transport's model and derivative information are shared.}
\label{tab:classical-information}
\end{table}

The degree-three panel uses 580 shared geometry exponentials, 2,520
additional paired boundary-kernel evaluations and 4,410 density evaluations;
its diagnostic is about 0.185 seconds, excluding imports and serialisation.
No speed ratio is inferred against joint incremental transport timings.
The independent 1,024-bit check reconstructs 180 moment enclosures,
144 coefficients, 108 remainders and all price-interval propagations.
The three failures are the same two-asset $\lambda=0.05$ cap-10,
cap-25 and ordinary calls; retained validated quadrature covers them.
The previous higher-degree checks remain unchanged. Strong classical
methods therefore still resolve all 36 prices with zero quantum calls.

I implement the published geometric bounds of \citet{caldana} from the
repository-labelled accepted manuscript, including its equation (27).
I write
\begin{equation}
\begin{aligned}
\sigma_G^2 &= \operatorname{Var}(\log G),\\
a_i &= \frac{\operatorname{Cov}(\log S_i,\log G)}{\sigma_G},\\
A_i &= w_i\E S_i.
\end{aligned}
\end{equation}
The lower-bound objective has derivative
$\phi(z)\{k-\sum_i A_ie^{a_i z-a_i^2/2}\}$.
Strict convexity handles the two mixed-sign covariance families; positive
basket weights do not imply a monotone conditional mean. The complete
root and endpoint proof, exact brackets and cap subtraction are retained
outside the main argument. This implemented enclosure meets no cent
allocation here, which does not defeat the published method's stronger
variants or its use as a control. Correctness and literature priority
remain separate from this numerical finding.

I test the observed admission coincidence beyond the original cohort
by multiplying each outer loading by $c\in\{1,3/2,2,3\}$ and recomputing
its risk-neutral drift. The common $\lambda$, weights, carry, contracts
and numerical widths remain fixed. This gives 144 contracts. Exact
scaling gives $d_i(c)=c^2d_i(1)$, $s_i(c)=cs_i(1)$ and
$R_G(c)=cR_G(1)$; every asset forward remains $100e^{0.01}$.
The supporting proof derives these identities and the unchanged pricing
conditions. For $V_R=\sum_iw_i\|s_i(1)\|^2>0$, Jensen also gives
$\E U_c^k\geq\{e^{c^2V_R/2}-1\}^k$. Thus the raw fixed-degree
Taylor radius can grow without the price range becoming unbounded.
This does not prove difficulty for all classical pricing methods.

The fixed panel admits 81 transfer allowances and 71 strengthened
Taylor-price intervals through degree three. Ten cases separate those
criteria; targeted degree five resolves six and leaves four. I intersect
all available lower orders and intrinsic payoff ranges before classifying
a price failure. This opponent strengthening follows the first executed
run and is recorded transparently; its panel and financial targets are
unchanged. Earlier executions and all adverse cases remain retained.
The four survivors are capped calls with $H=1$: $(n,\lambda,c)$ equals
$(2,0.4,3)$, $(4,0.05,3/2)$, $(4,0.2,3)$ and $(8,0.2,3)$.
Their transfer allowances are approximately $0.00342842$, $0.00270704$,
$0.00466658$ and $0.00163955$, rounded upwards, while their best
strengthened degree-five price radii are approximately $0.0279505$,
$0.0305001$, $0.0509043$ and $0.0127283$, also rounded upwards.
This establishes useful transfer certification beyond this declared
finite Taylor screen. It establishes neither complete quantum prices for
these stresses nor failure of stronger classical algorithms. The original
quantum resource caps are not extrapolated to altered supports or
normalisers. Figure~\ref{fig:dispersion-separation} shows the distinct
certificate behaviour and the effect of the stronger classical check.

% Self-contained scientific figure: standard xcolor and graphicx only.
\begin{figure}[!htb]%
\centering%
\begingroup%
\definecolor{QBPStressBlue}{HTML}{0072B2}%
\definecolor{QBPStressOrange}{HTML}{D55E00}%
\definecolor{QBPStressGrid}{HTML}{DDDDDD}%
\setlength{\unitlength}{1mm}%
\resizebox{\linewidth}{!}{%
\begin{picture}(160,75)(0,10)%
\fontsize{8}{10}\selectfont%
\color{black}%
\put(48.5,82){\makebox(0,0){(a) Eight-asset ordinary-call family}}%
\put(128.5,82){\makebox(0,0){(b) Complete stress panel}}%
\color{QBPStressGrid}\linethickness{0.18pt}%
\qbezier(17.0000,23.0000)(48.5000,23.0000)(80.0000,23.0000)%
\color{black}\put(15.000,23.0000){\makebox(0,0)[r]{$10^{-3}$}}%
\color{QBPStressGrid}\linethickness{0.18pt}%
\qbezier(17.0000,33.6000)(48.5000,33.6000)(80.0000,33.6000)%
\color{black}\put(15.000,33.6000){\makebox(0,0)[r]{$10^{-2}$}}%
\color{QBPStressGrid}\linethickness{0.18pt}%
\qbezier(17.0000,44.2000)(48.5000,44.2000)(80.0000,44.2000)%
\color{black}\put(15.000,44.2000){\makebox(0,0)[r]{$10^{-1}$}}%
\color{QBPStressGrid}\linethickness{0.18pt}%
\qbezier(17.0000,54.8000)(48.5000,54.8000)(80.0000,54.8000)%
\color{black}\put(15.000,54.8000){\makebox(0,0)[r]{$10^{0}$}}%
\color{QBPStressGrid}\linethickness{0.18pt}%
\qbezier(17.0000,65.4000)(48.5000,65.4000)(80.0000,65.4000)%
\color{black}\put(15.000,65.4000){\makebox(0,0)[r]{$10^{1}$}}%
\color{QBPStressGrid}\linethickness{0.18pt}%
\qbezier(17.0000,76.0000)(48.5000,76.0000)(80.0000,76.0000)%
\color{black}\put(15.000,76.0000){\makebox(0,0)[r]{$10^{2}$}}%
\color{black}\linethickness{0.4pt}%
\qbezier(17.0000,23.0000)(48.5000,23.0000)(80.0000,23.0000)%
\qbezier(17.0000,23.0000)(17.0000,49.5000)(17.0000,76.0000)%
\put(3,23){\rotatebox{90}{\makebox(53,0){Allowance/its allocation}}}%
\qbezier(21.1087,23.0000)(21.1087,22.5000)(21.1087,22.0000)%
\put(21.1087,20){\makebox(0,0){1}}%
\qbezier(34.8043,23.0000)(34.8043,22.5000)(34.8043,22.0000)%
\put(34.8043,20){\makebox(0,0){1.5}}%
\qbezier(48.5000,23.0000)(48.5000,22.5000)(48.5000,22.0000)%
\put(48.5000,20){\makebox(0,0){2}}%
\qbezier(75.8913,23.0000)(75.8913,22.5000)(75.8913,22.0000)%
\put(75.8913,20){\makebox(0,0){3}}%
\put(48.5,15){\makebox(0,0){Loading scale $c$}}%
\color{black}\linethickness{0.5pt}%
\qbezier(17.0000,54.8000)(17.9000,54.8000)(18.8000,54.8000)%
\qbezier(19.8000,54.8000)(20.7000,54.8000)(21.6000,54.8000)%
\qbezier(22.6000,54.8000)(23.5000,54.8000)(24.4000,54.8000)%
\qbezier(25.4000,54.8000)(26.3000,54.8000)(27.2000,54.8000)%
\qbezier(28.2000,54.8000)(29.1000,54.8000)(30.0000,54.8000)%
\qbezier(31.0000,54.8000)(31.9000,54.8000)(32.8000,54.8000)%
\qbezier(33.8000,54.8000)(34.7000,54.8000)(35.6000,54.8000)%
\qbezier(36.6000,54.8000)(37.5000,54.8000)(38.4000,54.8000)%
\qbezier(39.4000,54.8000)(40.3000,54.8000)(41.2000,54.8000)%
\qbezier(42.2000,54.8000)(43.1000,54.8000)(44.0000,54.8000)%
\qbezier(45.0000,54.8000)(45.9000,54.8000)(46.8000,54.8000)%
\qbezier(47.8000,54.8000)(48.7000,54.8000)(49.6000,54.8000)%
\qbezier(50.6000,54.8000)(51.5000,54.8000)(52.4000,54.8000)%
\qbezier(53.4000,54.8000)(54.3000,54.8000)(55.2000,54.8000)%
\qbezier(56.2000,54.8000)(57.1000,54.8000)(58.0000,54.8000)%
\qbezier(59.0000,54.8000)(59.9000,54.8000)(60.8000,54.8000)%
\qbezier(61.8000,54.8000)(62.7000,54.8000)(63.6000,54.8000)%
\qbezier(64.6000,54.8000)(65.5000,54.8000)(66.4000,54.8000)%
\qbezier(67.4000,54.8000)(68.3000,54.8000)(69.2000,54.8000)%
\qbezier(70.2000,54.8000)(71.1000,54.8000)(72.0000,54.8000)%
\qbezier(73.0000,54.8000)(73.9000,54.8000)(74.8000,54.8000)%
\qbezier(75.8000,54.8000)(76.7000,54.8000)(77.6000,54.8000)%
\qbezier(78.6000,54.8000)(79.3000,54.8000)(80.0000,54.8000)%
\color{QBPStressBlue}\linethickness{0.65pt}%
\qbezier(21.1087,38.0906)(27.9565,42.2325)(34.8043,46.3744)%
\qbezier(34.8043,46.3744)(41.6522,49.5662)(48.5000,52.7579)%
\qbezier(48.5000,52.7579)(62.1957,58.0177)(75.8913,63.2775)%
\put(21.1087,38.0906){\circle*{1.2}}%
\put(34.8043,46.3744){\circle*{1.2}}%
\put(48.5000,52.7579){\circle*{1.2}}%
\put(75.8913,63.2775){\circle*{1.2}}%
\color{QBPStressOrange}\linethickness{0.65pt}%
\qbezier(21.1087,26.1776)(27.9565,33.7994)(34.8043,41.4212)%
\qbezier(34.8043,41.4212)(41.6522,46.9407)(48.5000,52.4602)%
\qbezier(48.5000,52.4602)(62.1957,60.5910)(75.8913,68.7219)%
\put(20.5587,25.6276){\rule{1.1mm}{1.1mm}}%
\put(34.2543,40.8712){\rule{1.1mm}{1.1mm}}%
\put(47.9500,51.9102){\rule{1.1mm}{1.1mm}}%
\put(75.3413,68.1719){\rule{1.1mm}{1.1mm}}%
\color{QBPStressBlue}\put(21,70){\circle*{1.2}}%
\color{black}\put(23,70){\makebox(0,0)[l]{Grid transfer}}%
\color{QBPStressOrange}\put(20.45,65.45){\rule{1.1mm}{1.1mm}}%
\color{black}\put(23,66){\makebox(0,0)[l]{Taylor $\leq3$ price interval}}%
\color{QBPStressGrid}\linethickness{0.18pt}%
\qbezier(103.0000,23.0000)(128.5000,23.0000)(154.0000,23.0000)%
\color{black}\put(101.000,23.0000){\makebox(0,0)[r]{0}}%
\color{QBPStressGrid}\linethickness{0.18pt}%
\qbezier(103.0000,33.6000)(128.5000,33.6000)(154.0000,33.6000)%
\color{black}\put(101.000,33.6000){\makebox(0,0)[r]{1}}%
\color{QBPStressGrid}\linethickness{0.18pt}%
\qbezier(103.0000,44.2000)(128.5000,44.2000)(154.0000,44.2000)%
\color{black}\put(101.000,44.2000){\makebox(0,0)[r]{2}}%
\color{QBPStressGrid}\linethickness{0.18pt}%
\qbezier(103.0000,54.8000)(128.5000,54.8000)(154.0000,54.8000)%
\color{black}\put(101.000,54.8000){\makebox(0,0)[r]{3}}%
\color{QBPStressGrid}\linethickness{0.18pt}%
\qbezier(103.0000,65.4000)(128.5000,65.4000)(154.0000,65.4000)%
\color{black}\put(101.000,65.4000){\makebox(0,0)[r]{4}}%
\color{QBPStressGrid}\linethickness{0.18pt}%
\qbezier(103.0000,76.0000)(128.5000,76.0000)(154.0000,76.0000)%
\color{black}\put(101.000,76.0000){\makebox(0,0)[r]{5}}%
\color{black}\linethickness{0.4pt}%
\qbezier(103.0000,23.0000)(128.5000,23.0000)(154.0000,23.0000)%
\qbezier(103.0000,23.0000)(103.0000,49.5000)(103.0000,76.0000)%
\put(89,23){\rotatebox{90}{\makebox(53,0){Separating contracts}}}%
\color{black}\put(109.3750,20){\makebox(0,0){1}}%
\color{QBPStressBlue}\put(118.3250,23.0000){\rule{3.5mm}{31.8000mm}}%
\color{QBPStressOrange}\put(122.4250,23.0000){\rule{3.5mm}{10.6000mm}}%
\color{black}\put(122.1250,20){\makebox(0,0){1.5}}%
\color{QBPStressBlue}\put(131.0750,23.0000){\rule{3.5mm}{31.8000mm}}%
\color{black}\put(134.8750,20){\makebox(0,0){2}}%
\color{QBPStressBlue}\put(143.8250,23.0000){\rule{3.5mm}{42.4000mm}}%
\color{QBPStressOrange}\put(147.9250,23.0000){\rule{3.5mm}{31.8000mm}}%
\color{black}\put(147.6250,20){\makebox(0,0){3}}%
\put(128.5,15){\makebox(0,0){Loading scale $c$}}%
\color{QBPStressBlue}\put(106,71){\rule{2mm}{1mm}}%
\color{black}\put(110,71.5){\makebox(0,0)[l]{After Taylor $\leq3$}}%
\color{QBPStressOrange}\put(106,67){\rule{2mm}{1mm}}%
\color{black}\put(110,67.5){\makebox(0,0)[l]{After targeted $\leq5$}}%
\end{picture}%
}%

\endgroup%
\small\setlength{\abovecaptionskip}{6pt}%
\caption{Financial-transfer and price certificates under a controlled dispersion stress. Panel (a) uses the eight-asset ordinary-call family at $\lambda=0.2$, selected before outcomes. Transfer allowances are divided by $0.005$ and complete Taylor-price intersection radii by $0.01$; the dashed line marks admission at one. These are different error targets. Joined points guide the eye between tested scales. Panel (b) counts transfer-admitted contracts whose strengthened Taylor-price intervals fail, among all 36 contracts per scale. Degree five is evaluated only for provisional degree-three separators. Price intervals include lower-order intersections and intrinsic bounds. These are synthetic certified allowances, not observed errors, quantum prices or speed-ups.}
\label{fig:dispersion-separation}
\end{figure}

I next price the four cap-one candidates completely. A direct conditional
price certificate combines the bounded kernel with global log-basket
coordinate derivatives; Appendix~\ref{app:direct-capped-price} derives the
financial transfer and full confidence interval. I sum all 1,024 grid scores
for the two-asset case. For each remaining case, a fixed 100,000-draw
calculation uses the non-asymptotic inequality of \citet{hoeffding1963},
with statistical radius $1/150$ and failure allowance $1/3600$.
Actual 10-bit labels come from retained operating-system words; coverage
explicitly assumes independent uniform words. Every evaluated score is
enclosed and rounded outwards to 40-bit endpoints. The whole Gaussian
tail and retained mass are included. All four complete undiscounted price
radii are below $0.01$.

\begin{table}[ht]
\centering\small
\begin{tabular}{rrrlrr}
\toprule
Assets&$\lambda$&$c$&Classical method&Undiscounted interval&Selected $R_1$ T cap\\
\midrule
2&0.4&3&Exact sum&$[0.406517,0.406908]$&25,306,395,925\\
4&0.05&1.5&Conditional MC&$[0.479091,0.496631]$&90,433,832,750\\
4&0.2&3&Conditional MC&$[0.437322,0.452372]$&97,976,923,250\\
8&0.2&3&Conditional MC&$[0.442736,0.460092]$&216,055,259,050\\
\bottomrule
\end{tabular}
\caption{Complete classical prices and independently constructed quantum
resource caps for the four synthetic cap-one candidates. Classical interval
endpoints are rounded outwards to six decimals. The two-asset exact sum is
deterministic; each Monte Carlo price has failure allowance $1/3600$ under
the stated source premise. Quantum columns are ideal-logical resource plans,
not measured prices or running times, with the same per-price failure
budget and radius $0.01$. The informed positive first correction is cheapest
within this earlier sixteen-target quantum family. These additive clean
caps precede the cancellation-aware comparison in Table~\ref{tab:shared-threshold}. Exact results,
financial errors, all alternatives and source hashes remain retained.}
\label{tab:four-complete-prices}
\end{table}

This earlier quantum construction costs all sixteen method targets, both
coherent lookup architectures and every monetary obligation. One eight-asset boundary
baseline fails with deterministic allowance approximately $0.01673354$;
the declared scalar-address-width-14 repair passes at approximately
$0.00441931$. The original rejection is retained. The direct first-residual
certificate is supplied to the informed $R_1$ opponent by taking its minimum
with the older bound; no classical Monte Carlo answer is imported into that
prior. The earlier first-correction caps range from approximately 25.3 to
216.1 billion logical T gates. Raw second corrections require approximately
0.301 to 2.045 trillion, and the boundary construction is not the cheapest
in any of these four cases. These are conditional costs in the declared
baseline and repair policy, not globally optimised resource bounds.

Independent replays verify financial enclosures, complete confidence and
emitted resource charges. Full traces, rejected plans, scalar tables,
source versions and reconstruction records remain in the supporting
archive. The reconstruction protocol precedes fresh execution with earlier
expected outcomes known; no blind first-outcome registration is claimed.
All retained evidence is synthetic. The complete comparison resolves the
four candidates classically and supplies valid conditional quantum plans.
Figure~\ref{fig:dispersion-separation} identifies failures of the declared
Taylor screen, rather than prices unavailable to classical calculation.

\subsection{A certified factor-basis comparison}
\label{sec:factor-basis}

Orthogonal Gaussian reparameterisation is established in classical
derivative pricing \citep{imai2006}. For an exact orthogonal matrix $O$,
replacing $R$ by $RO$ preserves $RR^{\mathsf T}$, risk-neutral drifts and
the continuous asset law because $OY$ has the law of $Y$. It also preserves
$c_0,c_1$ and continuous $U$ moments. The truncated and rounded product
laws generally change. With full loading maxima $A_j(O)=\max_i|(RO)_{ij}|$
and $r_j(O)=(R_G O)_j$, the direct cap-one transfer becomes
\begin{align}
B_1(O)&=\sum_j\{A_j(O)+|r_j(O)|\},\\
B_2(O)&=\sum_j A_j(O)^2+\|R_G\|_2^2,\\
\beta(O)&=a^{p-1}\{b_1M_1B_1(O)+b_2(M_1+M_2)B_2(O)\}+(1-a^p).
\end{align}
Here $a$ is one-coordinate retained mass, $b_1$ is the sum of absolute
cell first moments and $b_2$ is half the summed second moment. The earlier
global derivative proof applies in the new coordinates; cell Taylor
expansion and tensor telescoping give this allowance. The full invariance
and finite-law proof is retained in the supporting derivation.

I register identity and each single plane with cosine $3/5$ and sine
$\pm4/5$, giving $1,7,7,43$ candidates across the four cases. Earlier
loading diagnostics were known. The exact minimum direct allowance selects
one candidate per case; I then rebuild every finite moment, support,
normaliser, arithmetic certificate and complete method cost for that basis
and identity. The direct bounds improve by about 25.25, 23.83 and 10.56
per cent in the three multi-factor cases, but complete selected T caps fall
by only 0.05, 0.16 and 2.07 per cent. The phase sizes remain unchanged.
One rotated stable-Gauss allocation fails and a boundary cost rises
substantially; both remain recorded. The two-asset identity is the control.
This is a restricted paired cost comparison, not cost optimisation over
all 58 bases. Every representation and certificate is available to the
classical opponents; their existing continuous price certificates remain
valid for the invariant target.

\subsection{A shared threshold pricing preparation}
\label{sec:shared-threshold}

I develop the existing positive kernel representation into a different
coherent first correction. It trades conditional-price evaluation for a
binary event; the financial law and analytic control remain unchanged.

\begin{proposition}[A positive coupled price correction]
\label{prop:shared-threshold}
For the finite nonnegative measure $\nu$ in
\eqref{eq:measure-kernel}, of mass $0<H<\infty$, let $X$ have law
$\nu/H$ and let $N$ be an independent standard normal variable. For
$\Delta\geq0$ and $\lambda>0$,
\begin{equation}
P_\nu(m+\Delta)-P_\nu(m)
=H\Pr\{m<X-\lambda N\leq m+\Delta\}.
\end{equation}
For a capped call, $X=\log(1+(H/K)U_0)$ with $U_0$ uniform on $[0,1]$.
Cash digitals and finite positive mixtures have the same representation;
the prepared threshold law must separately be certified.
\end{proposition}
Conditioning on the common threshold converts each endpoint to a Gaussian half-line probability. Subtracting nested events gives the positive correction while preserving its payoff normaliser. Appendix~\ref{app:proof-shared-threshold} gives the full proof.

The same auxiliary draw is essential for both endpoints. The score has
normaliser $H$ for every draw. In contrast,
$\Delta P'(m+t\Delta)$ is bounded by $\Delta_*M_1$; clipping that
integrand at $H$ can change its expectation. Its retained allowance screen
does not constitute a fully costed rival.

For the cap-one implementation I use 10-bit inner normal labels, four-bit
uniform strike labels and a six-bit strike-log table. Outer scalar
addresses remain at twelve bits. Let $\mathcal L$ be this finite set of
labels, with $|\mathcal L|=2^{10p+14}$. The clean reversible preparation is
\begin{align}
A_{\rm ev}|0\rangle
&=|\mathcal L|^{-1/2}\sum_{\ell\in\mathcal L}
 |\ell\rangle|b_\ell\rangle|0_{\rm work}\rangle,\\
b_\ell&=\mathbf1\{\widehat L_\ell<\widehat Z_\ell
                         \leq\widehat R_\ell\}.
\end{align}
Here $\widehat L_\ell,\widehat R_\ell,\widehat Z_\ell$ are the actual
represented lower, upper and threshold registers with a common exact
offset $J$. Removing $J$ gives the financial endpoints and threshold.
Two ordered comparisons and their inverses implement the interval flag;
all private work is cleaned before reflection. The represented upper
endpoint is the same lower register plus a nonnegative log-ratio, with
range guards excluding overflow. Canonical estimation uses this actual
flag and its actual reflected core.

The added approximation is explicit. With $M_S=16$, $M_N=1024$,
$a_N=\Phi(6)-\Phi(-6)$ and normal node displacement $\delta_N$, take
checked $\nu_0\geq\sup\phi$ and $\nu_1\geq\sup|\phi'|$ and put
\begin{align}
\beta_S/H&=\frac{(H/K)^2}{12M_S^2}
               (\nu_1/\lambda^2+\nu_0/\lambda),\\
D_N&=(1-a_N)/2+a_N/(2M_N)+\nu_0\delta_N,\\
\eta&=\min\{1,\beta_S/H+2D_N+
                    \nu_0(e_\ell+e_u+2e_z)/\lambda\}.
\end{align}
The $e$ terms bound actual endpoint and threshold arithmetic. A midpoint
Taylor bound proves $\beta_S$; a labelled quantile coupling bounds both
CDF conventions by $D_N$. Expanding and shrinking event endpoints against
the continuous Gaussian proves the arithmetic allowance, including empty
intervals and discrete ties. Finite atoms are not treated as continuous
points. The supporting proof gives every derivation and circuit invariant.
The outer transfer applies to the smooth conditional correction, and
the uniform $\eta$ then relates its nodal values to expected binary flags.
No derivative of a discontinuous flag is assumed. Only the outer mass is
decoded: inner normal truncation is already in $D_N$. If $\bar a_p$ bounds
the outer mass and $r_a$ is its radius, then on the certified estimation event
\begin{equation}
|\widehat V-V|\leq r_{c_0}+\beta_y+
             \bar a_pH(\eta+\varepsilon_{\rm AE})+r_aH.
\end{equation}
The complete plan selects $\varepsilon_{\rm AE}$ to meet radius $0.01$
with failure allowance $1/3600$.

I test the encoding at four already selected bases, with no repairs.
I then strengthen the numerical competitor by independently allowing
$b_A,b_G\in\{6,8,10,12\}$ for its two price tables, leaving other widths
fixed. This registered 64-configuration pool retains twelve rejected
budgets. Every admitted cost includes the same full financial and
confidence obligations. Its numerical minima use $b_A=b_G=10$.
Figure~\ref{fig:pricing-precision} exposes that numerical pool.

I also cost a direct common-normal rival. With
$f(x)=\min\{1,100(e^x-1)_+\}$ and $\Delta=\log T\geq0$, its score is
\begin{equation}
h_{m,\Delta}(N)=f(m+\Delta+\lambda N)-f(m+\lambda N)\in[0,1].
\end{equation}
Both payoffs use the same normal label. Exact monotone clamps restrict
exponential evaluation to $[0,1/64]$, since
$\log(101/100)<1/64$. The power-series inequality
$e^x\leq(1-x)^{-1}$ gives the rational output bound $64/63$.
I retain the initial broad-headroom, rational-range and common-width
studies. I then register independent geometric and arithmetic exponential
addresses $b_G,b_A\in\{8,9,10,11,12\}$ with normal widths
$\{8,10,12\}$. This gives 300 settings: 48 exact diagonal references and
252 fresh graphs, including the previously untested diagonal width nine.
All pointwise and smoothing budgets are admitted under the same conservative
certificate policy. Both lookup architectures and all outcomes remain retained.

Individually monotone approximations need not preserve their mutual order.
I retain the final projection in the actual probability score:
\begin{equation}
\begin{aligned}
\widehat h_y(n)&=\operatorname{clip}_{[0,1]}
 \bigl(\widehat f_A(\widehat X_A)-\widehat f_G(\widehat X_G)\bigr),\\
\operatorname{clip}_{[0,1]}(v)&=\min\{1,\max\{0,v\}\}.
\end{aligned}
\label{eq:projected-raw-score}
\end{equation}
An independently checked scalar-domain witness gives a negative represented
difference at equal arguments. The guarded positive subtraction masks this
modular difference before it becomes a financial value, and its inverse
cleans the workspace. Projection is $1$-Lipschitz and fixes the ordered
ideal payoff difference, so it preserves the charged approximation bound.
The full proof and exact witness are retained in the supporting appendix;
the witness does not assert occurrence in the four financial grids.

A pointwise payoff slope can overcharge shared endpoint errors. I integrate
the continuous normal before differentiating. For a positive threshold
measure $\mu$ of mass one, define
\begin{equation}
\begin{aligned}
P_{\lambda,\mu}(m)&=\int\Phi((m-a)/\lambda)\,\mu(da),\\
R_{\lambda,\mu}(m,\Delta)&=P_{\lambda,\mu}(m+\Delta)-P_{\lambda,\mu}(m),\\
\kappa_0&=\nu_0/\lambda,\qquad
C_\lambda(d)=\min\{\kappa_0,\nu_1d/\lambda^2\},\qquad d\geq0.
\end{aligned}
\label{eq:matched-common-shift-definition}
\end{equation}
Here $\nu_j$ bounds the unit Gaussian density derivative of order $j$.
Positivity gives $0\leq P'\leq\kappa_0$ and Gaussian smoothing gives
$|P''|\leq\nu_1/\lambda^2$. For a true gap
$0\leq\Delta\leq\Delta_*$, a common fixed shift of size $E_m$ and a
gap error $e_\Delta$ therefore satisfy
\begin{equation}
|R_{\lambda,\mu}(\widehat m,\widehat\Delta)
       -R_{\lambda,\mu}(m,\Delta)|
\leq C_\lambda(\Delta_*)E_m+\kappa_0e_\Delta.
\label{eq:matched-common-shift-stability}
\end{equation}
The critical step perturbs the gap at fixed $\widehat m$, then shifts
both endpoints at the true gap. A fixed strike, positive strike mixture
and the capped payoff derivative measure satisfy the same argument.
For raw payoffs, $f'\in[0,101]$ also makes the difference
$f(x+\widehat\Delta)-f(x)$ $101$-Lipschitz. With intrinsic exponential
errors $e_{\exp,G},e_{\exp,A}$, common inner scale error $E_s$, outer
shift error $E_m$ and inner-law CDF discrepancy $D_N$, let
$\eta_{\rm point}$ denote the retained valid pointwise conditional bound and
$\eta_{\rm previous}$ the minimum of previously valid conditional smoothing
bounds for the same graph and target. The new raw rule is
\begin{equation}
\begin{aligned}
\eta_{\rm raw}=\min\{\eta_{\rm point},\eta_{\rm previous},1,\,
&100(e_{\exp,G}+e_{\exp,A})+101E_s+2D_N\\
&+C_\lambda(\Delta_*)E_m+\kappa_0e_\Delta\}.
\end{aligned}
\label{eq:matched-common-shift-raw}
\end{equation}
The scaling remainder depends on the normal label and retains a pointwise
allowance. Treating it as a fixed shift under Gaussian smoothing is invalid.
The finite-to-continuous transfer precedes smoothing, retains atomic
multiplicities and charges normal rounding once. Projection cannot enlarge
the separate scalar or shared-argument errors.

I sharpen the positive-measure component bound without supplying a price
answer. Put $r_t(x)=\Phi(x+t/2)-\Phi(x-t/2)$ and define
\begin{equation}
\begin{aligned}
k(t)&=\sup_x|\phi(x+t)-\phi(x)|, &
\mathfrak a(t)&=2\Phi(t/2)-1,\\
\Omega(t,e)&=\sup_{x,\,|h|\leq e}|r_t(x+h)-r_t(x)|, &
t,e&\geq0.
\end{aligned}
\label{eq:sharp-component-definition-main}
\end{equation}
The sharp derivative coefficient is $k(d/\lambda)/\lambda$. The same
two-segment path yields the finite bound
\begin{equation}
|R_{\lambda,\mu}(\widehat m,\widehat\Delta)-R_{\lambda,\mu}(m,\Delta)|
\leq\Omega(\Delta_*/\lambda,E_m/\lambda)
                  +\mathfrak a(e_\Delta/\lambda).
\label{eq:sharp-component-path-main}
\end{equation}
Atomic measures attain each component extremum; their sum need not be
jointly sharp. Unique scalar roots certify these extrema outward.
The full proof also derives a stricter cap-specific envelope, with a
uniform gap certificate rather than assumed monotonicity. Interval checks
show derivative-component improvements of 7.86 and 16.47 per cent at
$d/\lambda=1$ and $1.5$. A cap-specific example improves that component
by at least 54.70 per cent. These are analytic certificate improvements,
not T-gate percentages, and are not selected in the registered resource
comparison below. Exact kernel-translation sensitivity has close published
antecedents \citep{duchi2012}; its consulted author text concerns an
$L^1$ gradient bound, whereas this extremum uses a positive measure and
an $L^\infty$ kernel difference. Sharpness in this class establishes no
literature priority.

I register the applicable refinement for all 348 raw, 64 numerical and
four threshold configurations, using each graph's actual rounding terms.
The numerical formula requires inactive input clamps and its exactly
represented unit normaliser; both premises are checked. The threshold
formula conditions on the strike label and separately sandwiches the
normal-dependent scaling remainder. Full proofs and all three formulae
are retained in the supporting appendix. All 416 certificates tighten,
but selected widths, phase resolutions, repetitions and complete canonical
caps are unchanged. The true gap obeys $\Delta\leq U\leq U_*$, giving
the shared price-independent choice $\Delta_*=U_*$ and
$C_\lambda(U_*)=\kappa_0$ in all four cases. A first-moment variation
refinement supplies no additional gain.

The raw minima use $b_G=b_A=9$; the numerical minima use width ten.
All earlier configurations and outcomes remain retained as references.

Under the fixed graph, word precision, normal widths, unary emitters and
canonical controller, an optimistic charge floor now excludes either
payoff table at widths fifteen through twenty-two, irrespective of the
other feasible width. Each side contributes four unary lookups in the
complete preparation, with selector charge proportional to $2^b+b$.
The supporting proof derives the exact count and excludes both zero-shot
prior decisions. Lower omitted widths, unequal widths thirteen/fourteen,
shared arithmetic and other estimators remain open. The common-information
amendment preserves these circuit profiles and the selected canonical costs.

I next address the published optimisation missing from additive template
charges. In operator order, let
$A=U^\dagger P U\mathsf H$ be the clean preparation, where $\mathsf H$
prepares the input labels, $U$ computes the guarded graph and $P$ XORs
its predicate into the flag. Since $U$ does not act on that flag,
the alternative $B=P U\mathsf H=UA$ has identical good probability.
The same actual zero-reflection tape gives conjugate Grover operators.
Its inverse restores zero work before every reflection, including the
borrowed ladder scratch. Dirty terminal work is measured and reset on
every allocated wire. The controlled tape has its own clean-work
contract; it is not assumed to implement the ideal reflection on
arbitrary dirty scratch. The supporting proof verifies both domains,
the conditional global minus and unchanged QPE distributions.

Writing $F_U=T(U)$, the globally cancellation-optimised clean QPE network
is the $B$ network followed by one terminal $U^\dagger$. For $R$ repetitions
and phase size $M$,
\begin{equation}
T_{\rm clean,opt}=T_B+RF_U,\qquad
T_{\rm clean,add}-T_B=R(2M-1)F_U.
\label{eq:cleanup-aware-complete-comparison}
\end{equation}
Thus the large additive-template reduction mainly reflects established
cancellation. Removing terminal cleanup saves only 0.097 to 0.195 per cent
against the stronger clean comparator across the twelve method/case minima.
Both variants are registered for every admitted configuration and both
architectures. No gate is assigned a free cost, and no fresh whole-circuit
emission is claimed: exact retained primitive profiles are transformed and
their full compositions independently checked.

\begin{table}[ht]
\centering\small
\begin{tabular}{rrrrrr}
\toprule
Assets&$\lambda$&$c$&Rival T cap&Event T cap&Saving (\%)\\
\midrule
2&0.4&3&11,630,271,225&10,467,507,700&10.00\\
4&0.05&1.5&44,122,469,300&41,861,923,075&5.12\\
4&0.2&3&46,988,300,100&44,709,494,725&4.85\\
8&0.2&3&104,186,984,850&101,750,688,975&2.34\\
\bottomrule
\end{tabular}
\caption{Complete canonical-QPE caps after the same cancellation-aware preparation is applied to every admitted method and architecture. Each rival is the cheaper deferred numerical or raw-payoff preparation. Raw payoff supplies the first, third and fourth comparator rows; numerical conditional pricing supplies the second. Every alternative has the same model, basis, analytic mean information, monetary radius $0.01$ and failure allowance $1/3600$. Savings are rounded to two decimals; exact costs and ratios are retained. These are sufficient logical implementation caps under synthetic contracts and finite policies, not global precision optima, quantum price observations or physical runtimes.}
\label{tab:shared-threshold}
\end{table}

Against the strongest equally optimised implemented rivals, the event
representation reduces complete caps by 2.34 to 10.00 per cent. The
selected phase sizes and 25 repetitions are unchanged. Independent checks
reconstruct 808 deferred profiles, 1,616 phase plans and 3,232 canonical
or half-confidence variant charges, retaining all twelve rejected budgets.
The financial certificates, priors and controller parameters are preserved
exactly by conjugacy. Earlier additive costs remain historical.
Compatible QSP evaluators, additional arithmetic optimisations and broader
estimator families remain separate comparisons.

The broader identity is not an efficient-state-preparation theorem for
every $\nu$, and does not cover negative measures, infinite payoff mass
or negative basket weights. A classical opponent integrates both
auxiliaries to recover the smooth conditional score $F(Y)$; explicitly,
\begin{equation}
\operatorname{Var}(H\mathbf1_{\rm event})
=\operatorname{Var}(F(Y))+\E[F(Y)(H-F(Y))].
\end{equation}
The unconditioned classical score therefore has no variance improvement.
For these cap-one targets, all three encodings have the same ideal
global probability $\E F/H$ and Grover angle before discretisation.
The encoder does not improve the generic query law. The complete classical prices remain valid, and no practical quantum
advantage or priority for the established threshold identity is inferred.

% Self-contained vector figure; xcolor and graphicx only.
% Self-contained native vector figure; xcolor and graphicx only.
% Source variant: deferred_B; all integer caps and exact ratios are retained in the backing JSON.
% Retained sixty-sample viridis RGB knots; linear interpolation in logarithmic colour position.
\begin{figure}[!htb]%
\centering\begingroup%
\definecolor{QBPRejected}{HTML}{E2E2E2}%
\definecolor{QBPInterval}{HTML}{D55E00}%
\setlength{\unitlength}{1mm}%
\noindent\makebox[\linewidth][c]{%
\resizebox{0.96\linewidth}{!}{%
\begin{picture}(162,84)(-2,4)%
\fontsize{7}{9}\selectfont%
\put(80,85){\makebox(0,0){Complete numerical $B$ T cap/numerical $B$ minimum}}%
\definecolor{QBPBar0}{rgb}{0.269944,0.014625,0.341379}%
\color{QBPBar0}\put(25.00000,77){\rule{1.85333mm}{3mm}}%
\definecolor{QBPBar1}{rgb}{0.274952,0.037752,0.364543}%
\color{QBPBar1}\put(26.83333,77){\rule{1.85333mm}{3mm}}%
\definecolor{QBPBar2}{rgb}{0.278791,0.062145,0.386592}%
\color{QBPBar2}\put(28.66667,77){\rule{1.85333mm}{3mm}}%
\definecolor{QBPBar3}{rgb}{0.281446,0.084320,0.407414}%
\color{QBPBar3}\put(30.50000,77){\rule{1.85333mm}{3mm}}%
\definecolor{QBPBar4}{rgb}{0.283091,0.110553,0.431554}%
\color{QBPBar4}\put(32.33333,77){\rule{1.85333mm}{3mm}}%
\definecolor{QBPBar5}{rgb}{0.283072,0.130895,0.449241}%
\color{QBPBar5}\put(34.16667,77){\rule{1.85333mm}{3mm}}%
\definecolor{QBPBar6}{rgb}{0.281887,0.150881,0.465405}%
\color{QBPBar6}\put(36.00000,77){\rule{1.85333mm}{3mm}}%
\definecolor{QBPBar7}{rgb}{0.279574,0.170599,0.479997}%
\color{QBPBar7}\put(37.83333,77){\rule{1.85333mm}{3mm}}%
\definecolor{QBPBar8}{rgb}{0.275191,0.194905,0.496005}%
\color{QBPBar8}\put(39.66667,77){\rule{1.85333mm}{3mm}}%
\definecolor{QBPBar9}{rgb}{0.270595,0.214069,0.507052}%
\color{QBPBar9}\put(41.50000,77){\rule{1.85333mm}{3mm}}%
\definecolor{QBPBar10}{rgb}{0.265145,0.232956,0.516599}%
\color{QBPBar10}\put(43.33333,77){\rule{1.85333mm}{3mm}}%
\definecolor{QBPBar11}{rgb}{0.257322,0.256130,0.526563}%
\color{QBPBar11}\put(45.16667,77){\rule{1.85333mm}{3mm}}%
\definecolor{QBPBar12}{rgb}{0.250425,0.274290,0.533103}%
\color{QBPBar12}\put(47.00000,77){\rule{1.85333mm}{3mm}}%
\definecolor{QBPBar13}{rgb}{0.243113,0.292092,0.538516}%
\color{QBPBar13}\put(48.83333,77){\rule{1.85333mm}{3mm}}%
\definecolor{QBPBar14}{rgb}{0.235526,0.309527,0.542944}%
\color{QBPBar14}\put(50.66667,77){\rule{1.85333mm}{3mm}}%
\definecolor{QBPBar15}{rgb}{0.225863,0.330805,0.547314}%
\color{QBPBar15}\put(52.50000,77){\rule{1.85333mm}{3mm}}%
\definecolor{QBPBar16}{rgb}{0.218130,0.347432,0.550038}%
\color{QBPBar16}\put(54.33333,77){\rule{1.85333mm}{3mm}}%
\definecolor{QBPBar17}{rgb}{0.210503,0.363727,0.552206}%
\color{QBPBar17}\put(56.16667,77){\rule{1.85333mm}{3mm}}%
\definecolor{QBPBar18}{rgb}{0.203063,0.379716,0.553925}%
\color{QBPBar18}\put(58.00000,77){\rule{1.85333mm}{3mm}}%
\definecolor{QBPBar19}{rgb}{0.194100,0.399323,0.555565}%
\color{QBPBar19}\put(59.83333,77){\rule{1.85333mm}{3mm}}%
\definecolor{QBPBar20}{rgb}{0.187231,0.414746,0.556547}%
\color{QBPBar20}\put(61.66667,77){\rule{1.85333mm}{3mm}}%
\definecolor{QBPBar21}{rgb}{0.180629,0.429975,0.557282}%
\color{QBPBar21}\put(63.50000,77){\rule{1.85333mm}{3mm}}%
\definecolor{QBPBar22}{rgb}{0.172719,0.448791,0.557885}%
\color{QBPBar22}\put(65.33333,77){\rule{1.85333mm}{3mm}}%
\definecolor{QBPBar23}{rgb}{0.166617,0.463708,0.558119}%
\color{QBPBar23}\put(67.16667,77){\rule{1.85333mm}{3mm}}%
\definecolor{QBPBar24}{rgb}{0.160665,0.478540,0.558115}%
\color{QBPBar24}\put(69.00000,77){\rule{1.85333mm}{3mm}}%
\definecolor{QBPBar25}{rgb}{0.154815,0.493313,0.557840}%
\color{QBPBar25}\put(70.83333,77){\rule{1.85333mm}{3mm}}%
\definecolor{QBPBar26}{rgb}{0.147607,0.511733,0.557049}%
\color{QBPBar26}\put(72.66667,77){\rule{1.85333mm}{3mm}}%
\definecolor{QBPBar27}{rgb}{0.141935,0.526453,0.555991}%
\color{QBPBar27}\put(74.50000,77){\rule{1.85333mm}{3mm}}%
\definecolor{QBPBar28}{rgb}{0.136408,0.541173,0.554483}%
\color{QBPBar28}\put(76.33333,77){\rule{1.85333mm}{3mm}}%
\definecolor{QBPBar29}{rgb}{0.131172,0.555899,0.552459}%
\color{QBPBar29}\put(78.16667,77){\rule{1.85333mm}{3mm}}%
\definecolor{QBPBar30}{rgb}{0.125394,0.574318,0.549086}%
\color{QBPBar30}\put(80.00000,77){\rule{1.85333mm}{3mm}}%
\definecolor{QBPBar31}{rgb}{0.121831,0.589055,0.545623}%
\color{QBPBar31}\put(81.83333,77){\rule{1.85333mm}{3mm}}%
\definecolor{QBPBar32}{rgb}{0.119738,0.603785,0.541400}%
\color{QBPBar32}\put(83.66667,77){\rule{1.85333mm}{3mm}}%
\definecolor{QBPBar33}{rgb}{0.119699,0.618490,0.536347}%
\color{QBPBar33}\put(85.50000,77){\rule{1.85333mm}{3mm}}%
\definecolor{QBPBar34}{rgb}{0.123444,0.636809,0.528763}%
\color{QBPBar34}\put(87.33333,77){\rule{1.85333mm}{3mm}}%
\definecolor{QBPBar35}{rgb}{0.130067,0.651384,0.521608}%
\color{QBPBar35}\put(89.16667,77){\rule{1.85333mm}{3mm}}%
\definecolor{QBPBar36}{rgb}{0.140210,0.665859,0.513427}%
\color{QBPBar36}\put(91.00000,77){\rule{1.85333mm}{3mm}}%
\definecolor{QBPBar37}{rgb}{0.157851,0.683765,0.501686}%
\color{QBPBar37}\put(92.83333,77){\rule{1.85333mm}{3mm}}%
\definecolor{QBPBar38}{rgb}{0.175707,0.697900,0.491033}%
\color{QBPBar38}\put(94.66667,77){\rule{1.85333mm}{3mm}}%
\definecolor{QBPBar39}{rgb}{0.196571,0.711827,0.479221}%
\color{QBPBar39}\put(96.50000,77){\rule{1.85333mm}{3mm}}%
\definecolor{QBPBar40}{rgb}{0.220124,0.725509,0.466226}%
\color{QBPBar40}\put(98.33333,77){\rule{1.85333mm}{3mm}}%
\definecolor{QBPBar41}{rgb}{0.252899,0.742211,0.448284}%
\color{QBPBar41}\put(100.16667,77){\rule{1.85333mm}{3mm}}%
\definecolor{QBPBar42}{rgb}{0.281477,0.755203,0.432552}%
\color{QBPBar42}\put(102.00000,77){\rule{1.85333mm}{3mm}}%
\definecolor{QBPBar43}{rgb}{0.311925,0.767822,0.415586}%
\color{QBPBar43}\put(103.83333,77){\rule{1.85333mm}{3mm}}%
\definecolor{QBPBar44}{rgb}{0.344074,0.780029,0.397381}%
\color{QBPBar44}\put(105.66667,77){\rule{1.85333mm}{3mm}}%
\definecolor{QBPBar45}{rgb}{0.386433,0.794644,0.372886}%
\color{QBPBar45}\put(107.50000,77){\rule{1.85333mm}{3mm}}%
\definecolor{QBPBar46}{rgb}{0.421908,0.805774,0.351910}%
\color{QBPBar46}\put(109.33333,77){\rule{1.85333mm}{3mm}}%
\definecolor{QBPBar47}{rgb}{0.458674,0.816363,0.329727}%
\color{QBPBar47}\put(111.16667,77){\rule{1.85333mm}{3mm}}%
\definecolor{QBPBar48}{rgb}{0.496615,0.826376,0.306377}%
\color{QBPBar48}\put(113.00000,77){\rule{1.85333mm}{3mm}}%
\definecolor{QBPBar49}{rgb}{0.545524,0.838039,0.275626}%
\color{QBPBar49}\put(114.83333,77){\rule{1.85333mm}{3mm}}%
\definecolor{QBPBar50}{rgb}{0.585678,0.846661,0.249897}%
\color{QBPBar50}\put(116.66667,77){\rule{1.85333mm}{3mm}}%
\definecolor{QBPBar51}{rgb}{0.626579,0.854645,0.223353}%
\color{QBPBar51}\put(118.50000,77){\rule{1.85333mm}{3mm}}%
\definecolor{QBPBar52}{rgb}{0.678489,0.863742,0.189503}%
\color{QBPBar52}\put(120.33333,77){\rule{1.85333mm}{3mm}}%
\definecolor{QBPBar53}{rgb}{0.720391,0.870350,0.162603}%
\color{QBPBar53}\put(122.16667,77){\rule{1.85333mm}{3mm}}%
\definecolor{QBPBar54}{rgb}{0.762373,0.876424,0.137064}%
\color{QBPBar54}\put(124.00000,77){\rule{1.85333mm}{3mm}}%
\definecolor{QBPBar55}{rgb}{0.804182,0.882046,0.114965}%
\color{QBPBar55}\put(125.83333,77){\rule{1.85333mm}{3mm}}%
\definecolor{QBPBar56}{rgb}{0.855810,0.888601,0.097452}%
\color{QBPBar56}\put(127.66667,77){\rule{1.85333mm}{3mm}}%
\definecolor{QBPBar57}{rgb}{0.896320,0.893616,0.096335}%
\color{QBPBar57}\put(129.50000,77){\rule{1.85333mm}{3mm}}%
\definecolor{QBPBar58}{rgb}{0.935904,0.898570,0.108131}%
\color{QBPBar58}\put(131.33333,77){\rule{1.85333mm}{3mm}}%
\definecolor{QBPBar59}{rgb}{0.974417,0.903590,0.130215}%
\color{QBPBar59}\put(133.16667,77){\rule{1.85333mm}{3mm}}%
\color{black}\put(25.00000,74){\makebox(0,0){1}}%
\color{black}\put(55.26868,74){\makebox(0,0){1.25}}%
\color{black}\put(80.00000,74){\makebox(0,0){1.5}}%
\color{black}\put(119.02312,74){\makebox(0,0){2}}%
\color{black}\put(22,67){\makebox(0,0){2 assets}}%
\put(22,62){\makebox(0,0){$\lambda=0.4,c=3$}}%
\definecolor{QBPCell0_6_6}{rgb}{0.742124,0.873494,0.149382}%
\color{QBPCell0_6_6}\put(8,29){\rule{7mm}{7mm}}%
\color{black}\put(11.5,32.5){\makebox(0,0){2.075}}%
\definecolor{QBPCell0_6_8}{rgb}{0.702182,0.867478,0.174293}%
\color{QBPCell0_6_8}\put(8,36){\rule{7mm}{7mm}}%
\color{black}\put(11.5,39.5){\makebox(0,0){2.049}}%
\definecolor{QBPCell0_6_10}{rgb}{0.694160,0.866213,0.179442}%
\color{QBPCell0_6_10}\put(8,43){\rule{7mm}{7mm}}%
\color{black}\put(11.5,46.5){\makebox(0,0){2.043}}%
\definecolor{QBPCell0_6_12}{rgb}{0.806928,0.882395,0.114034}%
\color{QBPCell0_6_12}\put(8,50){\rule{7mm}{7mm}}%
\color{black}\put(11.5,53.5){\makebox(0,0){2.119}}%
\definecolor{QBPCell0_8_6}{rgb}{0.275364,0.040369,0.366909}%
\color{QBPCell0_8_6}\put(15,29){\rule{7mm}{7mm}}%
\color{white}\put(18.5,32.5){\makebox(0,0){1.022}}%
\definecolor{QBPCell0_8_8}{rgb}{0.270652,0.017893,0.344652}%
\color{QBPCell0_8_8}\put(15,36){\rule{7mm}{7mm}}%
\color{white}\put(18.5,39.5){\makebox(0,0){1.009}}%
\definecolor{QBPCell0_8_10}{rgb}{0.269944,0.014625,0.341379}%
\color{QBPCell0_8_10}\put(15,43){\rule{7mm}{7mm}}%
\color{white}\put(18.5,46.5){\makebox(0,0){1.006}}%
\definecolor{QBPCell0_8_12}{rgb}{0.280555,0.076882,0.400430}%
\color{QBPCell0_8_12}\put(15,50){\rule{7mm}{7mm}}%
\color{white}\put(18.5,53.5){\makebox(0,0){1.044}}%
\definecolor{QBPCell0_10_6}{rgb}{0.273285,0.030054,0.356833}%
\color{QBPCell0_10_6}\put(22,29){\rule{7mm}{7mm}}%
\color{white}\put(25.5,32.5){\makebox(0,0){1.016}}%
\definecolor{QBPCell0_10_8}{rgb}{0.269944,0.014625,0.341379}%
\color{QBPCell0_10_8}\put(22,36){\rule{7mm}{7mm}}%
\color{white}\put(25.5,39.5){\makebox(0,0){1.003}}%
\definecolor{QBPCell0_10_10}{rgb}{0.269944,0.014625,0.341379}%
\color{QBPCell0_10_10}\put(22,43){\rule{7mm}{7mm}}%
\color{white}\put(25.5,46.5){\makebox(0,0){1.000}}%
\color{white}\linethickness{1.6pt}%
\put(22,43){\framebox(7,7){}}%
\color{black}\linethickness{0.8pt}%
\put(22,43){\framebox(7,7){}}%
\definecolor{QBPCell0_10_12}{rgb}{0.279411,0.067326,0.391457}%
\color{QBPCell0_10_12}\put(22,50){\rule{7mm}{7mm}}%
\color{white}\put(25.5,53.5){\makebox(0,0){1.038}}%
\definecolor{QBPCell0_12_6}{rgb}{0.280232,0.074181,0.397893}%
\color{QBPCell0_12_6}\put(29,29){\rule{7mm}{7mm}}%
\color{white}\put(32.5,32.5){\makebox(0,0){1.042}}%
\definecolor{QBPCell0_12_8}{rgb}{0.277238,0.052277,0.377672}%
\color{QBPCell0_12_8}\put(29,36){\rule{7mm}{7mm}}%
\color{white}\put(32.5,39.5){\makebox(0,0){1.029}}%
\definecolor{QBPCell0_12_10}{rgb}{0.276507,0.047632,0.373474}%
\color{QBPCell0_12_10}\put(29,43){\rule{7mm}{7mm}}%
\color{white}\put(32.5,46.5){\makebox(0,0){1.026}}%
\definecolor{QBPCell0_12_12}{rgb}{0.283090,0.111986,0.432800}%
\color{QBPCell0_12_12}\put(29,50){\rule{7mm}{7mm}}%
\color{white}\put(32.5,53.5){\makebox(0,0){1.064}}%
\color{black}\put(11.5,26){\makebox(0,0){6}}%
\put(6.5,32.5){\makebox(0,0)[r]{6}}%
\color{black}\put(18.5,26){\makebox(0,0){8}}%
\put(6.5,39.5){\makebox(0,0)[r]{8}}%
\color{black}\put(25.5,26){\makebox(0,0){10}}%
\put(6.5,46.5){\makebox(0,0)[r]{10}}%
\color{black}\put(32.5,26){\makebox(0,0){12}}%
\put(6.5,53.5){\makebox(0,0)[r]{12}}%
\put(22,22){\makebox(0,0){$b_A$}}%
\put(1,43){\rotatebox{90}{\makebox(0,0){$b_G$}}}%
\color{QBPInterval}\put(22,17){\makebox(0,0){Event $B$/min $B$=0.879}}%
\color{black}\put(61,67){\makebox(0,0){4 assets}}%
\put(61,62){\makebox(0,0){$\lambda=0.05,c=1.5$}}%
\color{QBPRejected}\put(47,29){\rule{7mm}{7mm}}%
\color{black}\put(50.5,32.5){\makebox(0,0){X}}%
\color{QBPRejected}\put(47,36){\rule{7mm}{7mm}}%
\color{black}\put(50.5,39.5){\makebox(0,0){X}}%
\color{QBPRejected}\put(47,43){\rule{7mm}{7mm}}%
\color{black}\put(50.5,46.5){\makebox(0,0){X}}%
\color{QBPRejected}\put(47,50){\rule{7mm}{7mm}}%
\color{black}\put(50.5,53.5){\makebox(0,0){X}}%
\definecolor{QBPCell1_8_6}{rgb}{0.670652,0.862369,0.194613}%
\color{QBPCell1_8_6}\put(54,29){\rule{7mm}{7mm}}%
\color{black}\put(57.5,32.5){\makebox(0,0){2.029}}%
\definecolor{QBPCell1_8_8}{rgb}{0.269944,0.014625,0.341379}%
\color{QBPCell1_8_8}\put(54,36){\rule{7mm}{7mm}}%
\color{white}\put(57.5,39.5){\makebox(0,0){1.007}}%
\definecolor{QBPCell1_8_10}{rgb}{0.269944,0.014625,0.341379}%
\color{QBPCell1_8_10}\put(54,43){\rule{7mm}{7mm}}%
\color{white}\put(57.5,46.5){\makebox(0,0){1.004}}%
\definecolor{QBPCell1_8_12}{rgb}{0.274233,0.034432,0.361217}%
\color{QBPCell1_8_12}\put(54,50){\rule{7mm}{7mm}}%
\color{white}\put(57.5,53.5){\makebox(0,0){1.019}}%
\definecolor{QBPCell1_10_6}{rgb}{0.271143,0.020162,0.346925}%
\color{QBPCell1_10_6}\put(61,29){\rule{7mm}{7mm}}%
\color{white}\put(64.5,32.5){\makebox(0,0){1.010}}%
\definecolor{QBPCell1_10_8}{rgb}{0.269944,0.014625,0.341379}%
\color{QBPCell1_10_8}\put(61,36){\rule{7mm}{7mm}}%
\color{white}\put(64.5,39.5){\makebox(0,0){1.003}}%
\definecolor{QBPCell1_10_10}{rgb}{0.269944,0.014625,0.341379}%
\color{QBPCell1_10_10}\put(61,43){\rule{7mm}{7mm}}%
\color{white}\put(64.5,46.5){\makebox(0,0){1.000}}%
\color{white}\linethickness{1.6pt}%
\put(61,43){\framebox(7,7){}}%
\color{black}\linethickness{0.8pt}%
\put(61,43){\framebox(7,7){}}%
\definecolor{QBPCell1_10_12}{rgb}{0.272796,0.027797,0.354572}%
\color{QBPCell1_10_12}\put(61,50){\rule{7mm}{7mm}}%
\color{white}\put(64.5,53.5){\makebox(0,0){1.015}}%
\definecolor{QBPCell1_12_6}{rgb}{0.274877,0.037405,0.364196}%
\color{QBPCell1_12_6}\put(68,29){\rule{7mm}{7mm}}%
\color{white}\put(71.5,32.5){\makebox(0,0){1.020}}%
\definecolor{QBPCell1_12_8}{rgb}{0.272206,0.025073,0.351844}%
\color{QBPCell1_12_8}\put(68,36){\rule{7mm}{7mm}}%
\color{white}\put(71.5,39.5){\makebox(0,0){1.013}}%
\definecolor{QBPCell1_12_10}{rgb}{0.271211,0.020477,0.347240}%
\color{QBPCell1_12_10}\put(68,43){\rule{7mm}{7mm}}%
\color{white}\put(71.5,46.5){\makebox(0,0){1.010}}%
\definecolor{QBPCell1_12_12}{rgb}{0.276149,0.045358,0.371418}%
\color{QBPCell1_12_12}\put(68,50){\rule{7mm}{7mm}}%
\color{white}\put(71.5,53.5){\makebox(0,0){1.025}}%
\color{black}\put(50.5,26){\makebox(0,0){6}}%
\put(45.5,32.5){\makebox(0,0)[r]{6}}%
\color{black}\put(57.5,26){\makebox(0,0){8}}%
\put(45.5,39.5){\makebox(0,0)[r]{8}}%
\color{black}\put(64.5,26){\makebox(0,0){10}}%
\put(45.5,46.5){\makebox(0,0)[r]{10}}%
\color{black}\put(71.5,26){\makebox(0,0){12}}%
\put(45.5,53.5){\makebox(0,0)[r]{12}}%
\put(61,22){\makebox(0,0){$b_A$}}%
\put(40,43){\rotatebox{90}{\makebox(0,0){$b_G$}}}%
\color{QBPInterval}\put(61,17){\makebox(0,0){Event $B$/min $B$=0.949}}%
\color{black}\put(100,67){\makebox(0,0){4 assets}}%
\put(100,62){\makebox(0,0){$\lambda=0.2,c=3$}}%
\color{QBPRejected}\put(86,29){\rule{7mm}{7mm}}%
\color{black}\put(89.5,32.5){\makebox(0,0){X}}%
\color{QBPRejected}\put(86,36){\rule{7mm}{7mm}}%
\color{black}\put(89.5,39.5){\makebox(0,0){X}}%
\color{QBPRejected}\put(86,43){\rule{7mm}{7mm}}%
\color{black}\put(89.5,46.5){\makebox(0,0){X}}%
\color{QBPRejected}\put(86,50){\rule{7mm}{7mm}}%
\color{black}\put(89.5,53.5){\makebox(0,0){X}}%
\definecolor{QBPCell2_8_6}{rgb}{0.272110,0.024627,0.351397}%
\color{QBPCell2_8_6}\put(93,29){\rule{7mm}{7mm}}%
\color{white}\put(96.5,32.5){\makebox(0,0){1.013}}%
\definecolor{QBPCell2_8_8}{rgb}{0.269944,0.014625,0.341379}%
\color{QBPCell2_8_8}\put(93,36){\rule{7mm}{7mm}}%
\color{white}\put(96.5,39.5){\makebox(0,0){1.006}}%
\definecolor{QBPCell2_8_10}{rgb}{0.269944,0.014625,0.341379}%
\color{QBPCell2_8_10}\put(93,43){\rule{7mm}{7mm}}%
\color{white}\put(96.5,46.5){\makebox(0,0){1.004}}%
\definecolor{QBPCell2_8_12}{rgb}{0.275608,0.041919,0.368310}%
\color{QBPCell2_8_12}\put(93,50){\rule{7mm}{7mm}}%
\color{white}\put(96.5,53.5){\makebox(0,0){1.023}}%
\definecolor{QBPCell2_10_6}{rgb}{0.270553,0.017439,0.344198}%
\color{QBPCell2_10_6}\put(100,29){\rule{7mm}{7mm}}%
\color{white}\put(103.5,32.5){\makebox(0,0){1.008}}%
\definecolor{QBPCell2_10_8}{rgb}{0.269944,0.014625,0.341379}%
\color{QBPCell2_10_8}\put(100,36){\rule{7mm}{7mm}}%
\color{white}\put(103.5,39.5){\makebox(0,0){1.002}}%
\definecolor{QBPCell2_10_10}{rgb}{0.269944,0.014625,0.341379}%
\color{QBPCell2_10_10}\put(100,43){\rule{7mm}{7mm}}%
\color{white}\put(103.5,46.5){\makebox(0,0){1.000}}%
\color{white}\linethickness{1.6pt}%
\put(100,43){\framebox(7,7){}}%
\color{black}\linethickness{0.8pt}%
\put(100,43){\framebox(7,7){}}%
\definecolor{QBPCell2_10_12}{rgb}{0.274267,0.034587,0.361373}%
\color{QBPCell2_10_12}\put(100,50){\rule{7mm}{7mm}}%
\color{white}\put(103.5,53.5){\makebox(0,0){1.019}}%
\definecolor{QBPCell2_12_6}{rgb}{0.273816,0.032507,0.359290}%
\color{QBPCell2_12_6}\put(107,29){\rule{7mm}{7mm}}%
\color{white}\put(110.5,32.5){\makebox(0,0){1.017}}%
\definecolor{QBPCell2_12_8}{rgb}{0.271304,0.020906,0.347670}%
\color{QBPCell2_12_8}\put(107,36){\rule{7mm}{7mm}}%
\color{white}\put(110.5,39.5){\makebox(0,0){1.010}}%
\definecolor{QBPCell2_12_10}{rgb}{0.270730,0.018256,0.345016}%
\color{QBPCell2_12_10}\put(107,43){\rule{7mm}{7mm}}%
\color{white}\put(110.5,46.5){\makebox(0,0){1.009}}%
\definecolor{QBPCell2_12_12}{rgb}{0.276903,0.050148,0.375748}%
\color{QBPCell2_12_12}\put(107,50){\rule{7mm}{7mm}}%
\color{white}\put(110.5,53.5){\makebox(0,0){1.028}}%
\color{black}\put(89.5,26){\makebox(0,0){6}}%
\put(84.5,32.5){\makebox(0,0)[r]{6}}%
\color{black}\put(96.5,26){\makebox(0,0){8}}%
\put(84.5,39.5){\makebox(0,0)[r]{8}}%
\color{black}\put(103.5,26){\makebox(0,0){10}}%
\put(84.5,46.5){\makebox(0,0)[r]{10}}%
\color{black}\put(110.5,26){\makebox(0,0){12}}%
\put(84.5,53.5){\makebox(0,0)[r]{12}}%
\put(100,22){\makebox(0,0){$b_A$}}%
\put(79,43){\rotatebox{90}{\makebox(0,0){$b_G$}}}%
\color{QBPInterval}\put(100,17){\makebox(0,0){Event $B$/min $B$=0.939}}%
\color{black}\put(139,67){\makebox(0,0){8 assets}}%
\put(139,62){\makebox(0,0){$\lambda=0.2,c=3$}}%
\color{QBPRejected}\put(125,29){\rule{7mm}{7mm}}%
\color{black}\put(128.5,32.5){\makebox(0,0){X}}%
\color{QBPRejected}\put(125,36){\rule{7mm}{7mm}}%
\color{black}\put(128.5,39.5){\makebox(0,0){X}}%
\color{QBPRejected}\put(125,43){\rule{7mm}{7mm}}%
\color{black}\put(128.5,46.5){\makebox(0,0){X}}%
\color{QBPRejected}\put(125,50){\rule{7mm}{7mm}}%
\color{black}\put(128.5,53.5){\makebox(0,0){X}}%
\definecolor{QBPCell3_8_6}{rgb}{0.642937,0.857512,0.212686}%
\color{QBPCell3_8_6}\put(132,29){\rule{7mm}{7mm}}%
\color{black}\put(135.5,32.5){\makebox(0,0){2.014}}%
\definecolor{QBPCell3_8_8}{rgb}{0.269944,0.014625,0.341379}%
\color{QBPCell3_8_8}\put(132,36){\rule{7mm}{7mm}}%
\color{white}\put(135.5,39.5){\makebox(0,0){1.003}}%
\definecolor{QBPCell3_8_10}{rgb}{0.269944,0.014625,0.341379}%
\color{QBPCell3_8_10}\put(132,43){\rule{7mm}{7mm}}%
\color{white}\put(135.5,46.5){\makebox(0,0){1.002}}%
\definecolor{QBPCell3_8_12}{rgb}{0.270329,0.016402,0.343159}%
\color{QBPCell3_8_12}\put(132,50){\rule{7mm}{7mm}}%
\color{white}\put(135.5,53.5){\makebox(0,0){1.008}}%
\definecolor{QBPCell3_10_6}{rgb}{0.635569,0.856220,0.217491}%
\color{QBPCell3_10_6}\put(139,29){\rule{7mm}{7mm}}%
\color{black}\put(142.5,32.5){\makebox(0,0){2.011}}%
\definecolor{QBPCell3_10_8}{rgb}{0.269944,0.014625,0.341379}%
\color{QBPCell3_10_8}\put(139,36){\rule{7mm}{7mm}}%
\color{white}\put(142.5,39.5){\makebox(0,0){1.001}}%
\definecolor{QBPCell3_10_10}{rgb}{0.269944,0.014625,0.341379}%
\color{QBPCell3_10_10}\put(139,43){\rule{7mm}{7mm}}%
\color{white}\put(142.5,46.5){\makebox(0,0){1.000}}%
\color{white}\linethickness{1.6pt}%
\put(139,43){\framebox(7,7){}}%
\color{black}\linethickness{0.8pt}%
\put(139,43){\framebox(7,7){}}%
\definecolor{QBPCell3_10_12}{rgb}{0.269944,0.014625,0.341379}%
\color{QBPCell3_10_12}\put(139,50){\rule{7mm}{7mm}}%
\color{white}\put(142.5,53.5){\makebox(0,0){1.006}}%
\definecolor{QBPCell3_12_6}{rgb}{0.651034,0.858931,0.207406}%
\color{QBPCell3_12_6}\put(146,29){\rule{7mm}{7mm}}%
\color{black}\put(149.5,32.5){\makebox(0,0){2.019}}%
\definecolor{QBPCell3_12_8}{rgb}{0.269944,0.014625,0.341379}%
\color{QBPCell3_12_8}\put(146,36){\rule{7mm}{7mm}}%
\color{white}\put(149.5,39.5){\makebox(0,0){1.005}}%
\definecolor{QBPCell3_12_10}{rgb}{0.269944,0.014625,0.341379}%
\color{QBPCell3_12_10}\put(146,43){\rule{7mm}{7mm}}%
\color{white}\put(149.5,46.5){\makebox(0,0){1.004}}%
\definecolor{QBPCell3_12_12}{rgb}{0.271109,0.020006,0.346768}%
\color{QBPCell3_12_12}\put(146,50){\rule{7mm}{7mm}}%
\color{white}\put(149.5,53.5){\makebox(0,0){1.010}}%
\color{black}\put(128.5,26){\makebox(0,0){6}}%
\put(123.5,32.5){\makebox(0,0)[r]{6}}%
\color{black}\put(135.5,26){\makebox(0,0){8}}%
\put(123.5,39.5){\makebox(0,0)[r]{8}}%
\color{black}\put(142.5,26){\makebox(0,0){10}}%
\put(123.5,46.5){\makebox(0,0)[r]{10}}%
\color{black}\put(149.5,26){\makebox(0,0){12}}%
\put(123.5,53.5){\makebox(0,0)[r]{12}}%
\put(139,22){\makebox(0,0){$b_A$}}%
\put(118,43){\rotatebox{90}{\makebox(0,0){$b_G$}}}%
\color{QBPInterval}\put(139,17){\makebox(0,0){Event $B$/min $B$=0.971}}%
\color{QBPRejected}\put(20,7){\rule{3mm}{3mm}}%
\color{black}\put(24,8.5){\makebox(0,0)[l]{X: rejected budget}}%
\linethickness{0.8pt}\put(73,7){\framebox(3,3){}}%
\put(77,8.5){\makebox(0,0)[l]{Exact numerical minimum}}%
\end{picture}%
}%
}\par%
\endgroup%
\small\setlength{\abovecaptionskip}{6pt}%
\caption{Complete T caps for the deferred-cleanup preparation $B$ over the 64 registered numerical-rival configurations. Panels give asset count, common volatility $\lambda$ and loading multiplier $c$. Columns and rows give the address widths $b_A$ and $b_G$ of $P(m_A)$ and $P(m_G)$; other widths and inputs are fixed. Each cell uses the cheaper complete $B$ cap across the two retained coherent architectures. Colours and labels show its cap/the numerical $B$ minimum in that panel, using the same shared logarithmic colour scale. Labels have three decimals; exact costs and ratios are retained. Grey crosses mark the same twelve rejected monetary budgets. Outlines identify the numerical minima, $b_A=b_G=10$. Orange footers compare the shared-threshold $B$ encoding with these numerical $B$ minima; Table~\ref{tab:shared-threshold} also includes the raw-payoff pool. The same established cancellation is applied to every representation, with each configuration retaining its financial certificate, information prior and phase plan. All alternatives use monetary radius $0.01$ and per-price failure allowance $1/3600$. The figure displays deterministic logical resource caps for synthetic models.}
\label{fig:pricing-precision}
\end{figure}


\subsection{Representation transfer across reflection samplers}

For a coherent preparation $A_i$ with a proved transported invariant,
let $a_i$ count its logical T gates
and $s_i$ those of its invariant input-zero reflection. The complete
adjoint has the same T cost. The good reflection is Clifford, and all
wire measurements and conditional Pauli resets retain their separately
charged zero-T logical cost. The result concerns additive retained
templates, with no cancellations across their boundaries.

\begin{lemma}[Matched-degree preparation dominance]
\label{lem:representation-sampler-dominance}
At every positive integer degree, the parity-specific reflection sampler
has complete per-shot T cost
\begin{equation}
t_i(d)=\begin{cases}
da_i+(d-1)s_i/2,&d\text{ odd},\\
da_i+(d/2-1)s_i,&d\text{ even}.
\end{cases}
\label{eq:representation-parity-cost}
\end{equation}
Candidate $c$ has no larger cost at every degree if and only if
\begin{equation}
a_c\leq a_b,\qquad2a_c+s_c\leq2a_b+s_b.
\label{eq:representation-all-degree-condition}
\end{equation}
For a rival with $a_b>0$, the sharp uniform cost ratio satisfies
\begin{equation}
\frac{t_c(d)}{t_b(d)}\leq
\max\left\{\frac{a_c}{a_b},\frac{2a_c+s_c}{2a_b+s_b}\right\}.
\label{eq:representation-all-degree-ratio}
\end{equation}
\end{lemma}
This is an implementation corollary of the published reflection samplers \citep{suzuki2020,rallfuller2023}. The structural step is that odd costs are $a_i+k(2a_i+s_i)$ and even
costs are $2a_i+(k-1)(2a_i+s_i)$. Two affine inequalities therefore
control every degree, including degrees outside the sampled panel.
Appendix~\ref{app:reflection-dominance} proves necessity, sufficiency,
the sharp ratio and the conditional same-history transfer.

The cancellation-aware threshold preparation satisfies both inequalities strictly
against each selected raw and numerical rival in all four cases. Its
per-shot saving therefore extends to the retained reflection-sampling
family. This is a project-specific resource consequence of the checked
pricing representation. Different rounding certificates permit different
probabilities, tolerances and adaptive histories. Identical implemented
probabilities and controller information are additional premises for
transferring the saving to entire adaptive histories.

\subsection{A matched bounded Chebyshev comparison}

I give each of the three cancellation-aware selected preparations the same declared
Chebyshev controller: at most twelve fresh 100-shot looks, degrees through
2,047, half the failure allowance for the pilot and half for a fresh
canonical fallback. Its complete protected cap is
$T_F+\lfloor T_F/10\rfloor$. A separately registered full-pool comparison
rebuilds all 808 admitted method/architecture fallbacks, preserving the
twelve rejected numerical budgets. Every fallback keeps its phase size,
Fourier tape and per-repetition charge, while repetitions increase from
25 to 27. Thus its protected-cap function is
\begin{equation}
\mathcal C(T_0)=\frac{27}{25}T_0+
                    \left\lfloor\frac{27}{250}T_0\right\rfloor,
\label{eq:matched-poolwise-protected-cap}
\end{equation}
which is strictly increasing on the admitted canonical caps. The same
precision minima therefore minimise these certified protected caps across
the full pool. Their relative ordering survives this controller change,
although each protected cap is approximately $1.188$ times its own
canonical cap. The modern branch therefore improves no worst-history
guarantee against canonical QPE here.

The diagnostic uses five stipulated probabilities in the common certified
ideal-mean interval and twenty paired seeds, giving 400 paired histories
and 1,200 alternative histories. Shared words remove method-relative
mesh and seed differences. They do not assert equality of the actual
rounded financial flags. Of these alternatives, 254 pilots certify over
1,968 confidence looks. Every median exceeds its own canonical cap.
The numerical competitor has the lower descriptive mean in every case,
even though its matched-degree shots cost more. Its distinct financial
allowances change its admissible degrees and stopping decisions.
Table~\ref{tab:matched-estimator-transfer} displays this limitation beside
the positive resource transfer. The supporting proof derives conditional
coverage, monetary matching, pairing, discarded-pilot charges and the
exact break-even criterion. Mesh means identify no expected cost at an
unknown financial flag. Unchanged published ChebAE heuristics, other
estimators and optimised upstream evaluators remain separate comparisons.

\begin{table}[ht]
\centering\small
\begin{tabular}{rrrrrr}
\toprule
Assets&$\lambda$&$c$&Shot floor (\%)&Protected saving (\%)&Mesh mean ratio\\
\midrule
2&0.4&3&9.98&10.00&1.088\\
4&0.05&1.5&5.11&5.12&1.025\\
4&0.2&3&4.84&4.85&1.088\\
8&0.2&3&2.33&2.34&1.054\\
\bottomrule
\end{tabular}
\caption{Cancellation-aware estimator transfer and its limit. The shot floor
bounds threshold savings against both selected rivals at every positive
integer degree and is rounded down. Protected savings use the full-pool
minimum certified caps and the cheaper equally optimised rival. The final
column gives threshold/numerical mean charge on the same stipulated mesh
and paired seeds; ratios above one favour numerical pricing. These are
synthetic resource comparisons, not expected-runtime or price estimates.}
\label{tab:matched-estimator-transfer}
\end{table}


The transport theorem applies beyond one payoff and asset count. It
certifies Gaussian integration whenever residual derivatives have finite
exponential-polynomial envelopes, including conditioned classical grids
and coherent financial preparation. For these baskets, moments through
degree four suffice for the second-residual transfer. Their construction
is polynomial in the number of exponential terms, with explicit scalar
setup and tensor-factor products; the error itself is not dimension-free.
The abstract finite-transform reduction shows that scalar factorisation
and Cauchy closure follow from established positive-measure tools. The
residual derivative calculation is a necessary financial specialisation,
with practical bounds, rather than a new general integration principle.
The checked argument supplies a reusable financial certificate and
reconstructed logical resource evidence. It establishes neither
superiority over classical pricing nor global precision optimality.
Compatible estimator comparisons beyond this bounded family and the four
stressed contracts, and broader market validation, remain necessary for
broader claims. The
supporting index maps complete derivations, source versions, failures and
byte-exact earlier archives and semantically reconstructible new evidence
without expanding the main argument.

\appendix
\section{Detailed proofs and standard derivations}
\label{app:main-derivations}

\subsection{Proof and scope of the reflection-sampler transfer}
\label{app:reflection-dominance}

\begin{proof}
Put $Q_i=-A_iS_{0,i}A_i^\dagger S_g$. The unconditional global minus
can be omitted because it changes a terminal shot only by a global
phase. It cannot be omitted from a controlled-QPE comparison. For
$d=2k+1$, the chronological sampler prepares $A_i$ and then applies
$Q_i$ exactly $k$ times. It uses $k+1$ preparations, $k$ adjoints,
$k$ zero reflections and $k$ Clifford good reflections. For $d=2k$
with $k\geq1$, it prepares $A_i$, applies $Q_i^{k-1}$, applies $S_g$
and then $A_i^\dagger$, observing the invariant input-zero event. It
uses $k$ preparations, $k$ adjoints, $k-1$ zero reflections and $k$
Clifford good reflections. Complete measurement and reset are appended
to both constructions. Their declared $T$ charge is zero, so these
counts prove \eqref{eq:representation-parity-cost}.

Let $q_i=2a_i+s_i$, $\delta a=a_c-a_b$ and $\delta q=q_c-q_b$.
The odd costs are $a_i+kq_i$ for $k\geq0$ and the even costs are
$2a_i+(k-1)q_i$ for $k\geq1$. If $\delta a\leq0$ and
$\delta q\leq0$, subtracting these formulae gives nonpositive
differences in both parities. Conversely, degree one requires
$\delta a\leq0$. If $\delta q>0$, the odd difference
$\delta a+k\delta q$ is positive at a sufficiently large integer
$k$, contradicting all-degree dominance. This proves necessity and
sufficiency. At $d=1$ strict dominance requires $\delta a<0$; together
with $\delta q\leq0$ this is also sufficient in both parities.

For odd degrees the ratio of costs is
$(a_c+kq_c)/(a_b+kq_b)$. It is a convex combination of $a_c/a_b$
and $q_c/q_b$, with nonnegative weights
$a_b/(a_b+kq_b)$ and $kq_b/(a_b+kq_b)$. For even degrees replace
$a_i$ by $2a_i$ and $k$ by $k-1$; the two endpoint ratios are
unchanged. A convex combination is bounded by the larger endpoint.
Degrees one and two attain $a_c/a_b$. Along either parity the ratio
tends to $q_c/q_b$ as its degree tends to infinity, proving the final
uniform sharpness statement.
\end{proof}

\begin{corollary}[Same-history adaptive sampler comparison]
\label{cor:representation-same-history}
Suppose two clean representations have the same implemented good
probability $p$, and an estimator uses the same initial prior, tolerance,
confidence allocation and deterministic history-dependent degree/batch
policy for both. Under ideal logical operation their complete reflection
shots can be coupled to have identical observed Bernoulli outcomes,
chosen degrees, batch sizes and stopping decisions. If
\eqref{eq:representation-all-degree-condition} holds, every coupled
history has no larger candidate reflection-shot $T$ cost. If
\eqref{eq:representation-all-degree-ratio} has right side $r$, then
the total candidate reflection-shot cost on every history is at most
$r$ times the corresponding rival reflection-shot cost. Consequently
the same comparison holds for expectations whenever they exist.
\end{corollary}
\begin{proof}
The two-dimensional reflection algebra gives the same
Chebyshev Bernoulli parameter $q_d(p)$ for both clean representations
at each degree. This includes $p=0,1$ by direct endpoint evaluation,
without normalising a zero-probability basis component. At each shot,
use a fresh common uniform random variable $V$ to define both observed
bits as $\mathbf1_{\{V\leq q_d(p)\}}$. Inductively, identical preceding
histories imply the same next degree and batch size, which imply
identical next observation under this coupling. The deterministic
policy then makes the same next decision. Apply the pointwise
per-shot inequality to each degree in the common history and sum its
nonnegative costs. Taking expectation preserves the inequality.
For unbounded stopping histories, apply the same argument to finite
prefixes and then use monotone convergence. Any additional fallback
or non-reflection primitive requires its own cost comparison; it is not
included by this corollary merely because it belongs to the estimator.
\end{proof}

\subsection{Proof of Theorem~\ref{thm:norm}}
\label{app:proof-norm}

\begin{proof}
The identity $\phi^{(j)}=(-1)^j\operatorname{He}_j\phi$ follows by
induction from $\phi'=-z\phi$ and the recurrence.  Its boundedness and
finite measure mass justify dominated differentiation under $\nu$.
If $D=\partial_m$, differentiating $t^{-r}Q(D)P(m+\log t)$ gives
$t^{-r-1}(D-r)Q(D)P(m+\log t)$.  Induction therefore yields
$g_m^{(r)}=t^{-r}\sum_\ell s(r,\ell)P_\nu^{(\ell)}$ and then the
integral in \eqref{eq:norm-derivative}.  Positivity of $\nu$ gives the
bound.  Since the leading term of $A_{r,\lambda}$ is
$(-1)^{r-1}\lambda^{-r}z^{r-1}$, the polynomial is nonzero.
Continuity and Gaussian decay give a maximiser $z_*$.  Take
$t=1$, $m=x_*+\lambda z_*$ and $\nu=H\delta_{x_*}$.
The absolute derivative is exactly $HD_r$, proving uniform sharpness.
\end{proof}

\subsection{Proof of Theorem~\ref{thm:finite-price}}
\label{app:proof-finite-price}

\begin{proof}
The integral Taylor remainder is
$k!^{-1}\int_1^T(T-t)^kg_m^{(k+1)}(t)dt$.
Since $T\geq1$, Theorem~\ref{thm:norm} bounds its absolute value by
$H\overline D_{k+1}U^{k+1}/(k+1)!$.
Every Gaussian exponential has all finite moments; finite sums preserve
this property, so the remainder is integrable.  Taking expectations gives
\eqref{eq:finite-certificate}.  The signed-Stirling derivative formula,
\eqref{eq:T-multinomial} and \eqref{eq:gaussian-tilt} evaluate all its
coefficients.  The finite collection of exponential and scalar normal
functions is computable at a positive scale, and outward arithmetic
preserves their enclosures.
\end{proof}

\subsection{Proof of Theorem~\ref{thm:boundary}}
\label{app:proof-boundary}

\begin{proof}
Linear growth gives an integrable payoff. For a single call, put
$x=m-\log(K_j/S_{\rm ref})$. Completing the Gaussian square gives
$P_j=K_j\{e^{x+\lambda^2/2}\Phi((x+\lambda^2)/\lambda)
-\Phi(x/\lambda)\}$.
The identity $e^{x+\lambda^2/2}\phi((x+\lambda^2)/\lambda)=\phi(x/\lambda)$
gives $P_j'=K_je^{x+\lambda^2/2}\Phi((x+\lambda^2)/\lambda)$ and
$P_j''-P_j'=K_j\phi(x/\lambda)/\lambda$.
Gaussian density times an exponential and a polynomial dominates every
required derivative on a compact $m$ interval, justifying differentiation.
The affine terms are killed by $D(D-1)$; finite linearity proves
\eqref{eq:boundary-kernel}. The chain rule gives
$g_m''(t)=t^{-2}C_2(m+\log t)$ for $g_m(t)=P_F(m+\log t)$.
Integrating $g_m''$ twice proves \eqref{eq:beta-identity}.
Its measure and $U^2$ are nonnegative, and $t\geq1$, proving both bounds.
\end{proof}

\subsection{Proof of Corollary~\ref{cor:convex-screen}}
\label{app:proof-convex-screen}

\begin{proof}
Theorem~\ref{thm:boundary} bounds the exact residual between zero and
$M_2U^2/2$ pointwise on the full Gaussian law. All lognormal ratio moments
are finite, so taking expectations proves both endpoints. Gaussian tilts
evaluate $c_0,c_1$ and the continuous moment in finite sums. Outward
coefficient errors enlarge the endpoints when they are implemented.
\end{proof}

\subsection{Proof of Proposition~\ref{prop:variance}}
\label{app:proof-variance}

\begin{proof}
Equation~\eqref{eq:beta-identity} gives the conditional mean. The finite
support makes second moments finite. Subtract the conditional mean and
expand the square; the mixed term has conditional expectation zero.
\end{proof}

\subsection{Proof of Theorem~\ref{thm:error}}
\label{app:proof-error}

\begin{proof}
Insert the exact coefficient, mass-weighted finite financial correction,
finite auxiliary expectation and implemented expectation in succession.
Their deterministic discrepancies are bounded by the stated terms.
The decoder is $B(2p-1)$ or $Bp$, giving $dB\eta_p$.
Write the external-factor discrepancy against the implemented bounded
score, giving $e_FB$. The triangle inequality proves the bound on the
estimator success event. Exact integer signed-part recombination preserves
the centred observable, so its deterministic bound is the centred one.
\end{proof}

\subsection{Proof of Proposition~\ref{prop:shared-threshold}}
\label{app:proof-shared-threshold}

\begin{proof}
Conditioning on $X=x$ gives
$\Pr\{x-\lambda N\leq m\}=\Phi((m-x)/\lambda)$.
Integrating against $\nu/H$ and subtracting the two nested half-line
events proves the identity. Gaussian convolution removes continuous
endpoint atoms. The cap law follows from the threshold payoff integral
and the change of variables $x=\log(1+(H/K)u)$.
\end{proof}

\subsection{Equivalent pricing measure and admissible trading}
\label{app:measure-change}

The strictly positive process
$L_t=\exp(-\theta\cdot W_t-\|\theta\|^2t/2)$ is a martingale of mean one:
the assertion follows by conditioning on independent Gaussian increments
and evaluating their exponential moment.  It defines an equivalent
probability $Q$ on the finite horizon.  Multiplying an increment's normal
density by its exponential factor shifts its mean by
$-\theta\Delta t$, preserving covariance and independence.  Therefore
$W_t^Q=W_t+\theta t$ is Brownian under $Q$, and the stock drift becomes
$r_f-q_i$.

The discounted cumulative gain
\begin{equation}
 e^{-r_ft}S_i(t)+\int_0^t e^{-r_fu}q_iS_i(u)\,du
\end{equation}
has zero drift under $Q$.  The explicit GBM solution has finite second
moments on $[0,\tau]$, so its stochastic integral is a true martingale.
For a self-financing admissible zero-cost strategy whose discounted wealth
is bounded below, localisation followed by conditional Fatou's lemma makes
the discounted wealth a supermartingale.  Nonnegative terminal wealth which
is positive with positive probability would contradict its nonpositive
expectation.  This establishes absence of arbitrage in the stated
admissibility class.  Full row rank is sufficient, rather than necessary.
If Brownian risk is unspanned, I use this stipulated $Q$; the quantum
estimator does not select a new pricing measure.

\subsection{Common-factor covariance decomposition}
\label{app:common-factor}

The equal common Gaussian loading is available for every positive-definite
maturity log-covariance matrix $C$. Let
$\lambda^2=(\mathbf1^{\mathsf T}C^{-1}\mathbf1)^{-1}$.
Cauchy--Schwarz applied to $C^{1/2}v$ and $C^{-1/2}\mathbf1$ gives
\begin{equation}
 v^{\mathsf T}Cv\geq
 \lambda^2(\mathbf1^{\mathsf T}v)^2,
\end{equation}
so $C-\lambda^2\mathbf1\mathbf1^{\mathsf T}$ is positive semidefinite.
It kills $C^{-1}\mathbf1$ and has rank $n-1$, because a rank-one
subtraction from positive-definite $C$ loses exactly that one direction.
A factor $RR^{\mathsf T}$ of this residual matrix and an independent
common $Z$ reproduce $C$ exactly. The original drifts use its unchanged
diagonal variances. This established Gaussian factorisation shows the
analytical model coverage; extracting a common factor is not itself a
new asset law or a proof of a favourable numerical cost.

\subsection{Gaussian tilt and differentiated conditioning}
\label{app:gaussian-tilt-detail}

For \eqref{eq:gaussian-tilt}, use
$e^{v\cdot y-\|y\|^2/2}=e^{\|v\|^2/2}e^{-\|y-v\|^2/2}$ and shift the
integration variable.  Independent Gaussian convolution changes the common
scale to $\sigma$.  Bounded differentiated kernels justify passing each
finite derivative through the expectation.  This remains valid at $R_G=0$.

\subsection{Closed-kernel differentiation and independent verification}
\label{app:kernel-differentiation}

With the variables in \eqref{eq:closed-kernel}, the density identities
$e^{m+\sigma^2/2}\phi(z_0)=\phi(x_0)$ and
$e^{m+\sigma^2/2}\phi(z_1)=(1+H/K)\phi(x_1)$ cancel the boundary terms
in differentiation.  They give
$P'_{\sigma,H}=Ke^{m+\sigma^2/2}[\Phi(z_0)-\Phi(z_1)]$;
higher derivatives follow by the product rule and Hermite recurrence.
The retained audit compares these evaluations with differentiated series
and the strike-integral representation.  A formal closed formula is not
used as its own verification.

\subsection{Separated Fourier and Hermite comparators}
\label{app:fourier-comparators}

For comparison with the earlier separated bounds, Gaussian Fourier
inversion gives
$\mathcal B_j=(2\pi)^{-1}\int_\R |u|^je^{-u^2/2}du$ and
$\mathcal H_j\leq\mathcal B_j$.
Direct integration gives $\mathcal B_0=(2\pi)^{-1/2}$,
$\mathcal B_1=\pi^{-1}$ and
integration by parts gives $\mathcal B_j=(j-1)\mathcal B_{j-2}$.
Thus the individual-norm triangle and Fourier triangle are distinct
competent comparators.  Retaining both separates a gain from cancellation
from a gain obtained by tightening each derivative separately.

\subsection{Canonical phase-grid error derivation}
\label{app:canonical-sketch}

Equation~\eqref{eq:canonical-error} follows from the two Grover eigenphases and the nearest phase-grid
outcomes: writing $p=\sin^2\theta$, the relevant phase error is at most
$\pi/M$, and Taylor expansion of $\sin^2$ bounds the first-order term by
$2\sqrt{p(1-p)}\pi/M$ and the second-order term by $\pi^2/M^2$.
The probability of the nearest-grid event is bounded below by the retained
canonical phase-estimation theorem.  The exact finite confidence calculation
and phase-tape verification remain in the supporting derivations.

\subsection{Fourier synthesis and retained tape diagnostics}
\label{app:fourier-diagnostics}

The inverse Fourier transform has explicit Clifford and T sequences,
based on certified single-qubit rotation approximations
\citep{rossselinger2016}.  Scalar phases are retained until they are
physically irrelevant to the final measurement.  For phase widths through
64, distant unsupported controlled phases are replaced by identity with a
charged operator-norm error; retained approximations and omitted phases
are summed by a telescoping product bound.
The complete largest-width bound is about $0.0022013$, below the stipulated
$1/64$ implementation limit.  The 64-bit tape has 167,223 T gates; the earlier 32-bit
profile is preserved at 68,823 T gates.  The extension charges omitted
distant phases rather than silently assigning them exact rotations.
This is an implementation-error statement, not an execution on a quantum
device.


\section{Complete root isolation for the joint norm}
\label{app:roots}

This appendix gives the finite certificate behind
Theorem~\ref{thm:norm}.  Write $A=A_{r,\lambda}$ and $Q=A'-zA$.
Since $(A\phi)'=Q\phi$, every nonzero absolute maximum is at a real root
of $Q$.  Zeros of $A$ and limits at infinity contribute zero.
For rational $\lambda$, all coefficients are rational and $\deg Q=r$.
Replace $Q$ by its monic square-free part $q=Q/\gcd(Q,Q')$; this keeps
every distinct root.

Start a remainder sequence with $q_0=q$, $q_1=q'$ and define
$q_{j+1}$ as minus the polynomial remainder of $q_{j-1}$ by $q_j$.
Positive rescalings do not change signs.  The sequence ends in a nonzero
constant.  Consecutive terms cannot vanish together, since the recurrence
would propagate that zero to the last constant.  At an internal zero the
neighbouring terms have opposite signs, so their adjacent comparisons
contribute one sign variation on both sides.  At a root of $q_0$, its
nonzero derivative makes the variation count decrease by one.
Thus, omitting zero values, $V(a)-V(b)$ counts roots in $(a,b)$ when
$a,b$ are nonroots.

Every root lies within
$1+\max_{j<r}|a_j/a_r|$ for the coefficients of $q$.
Outside that radius the absolute sum of its lower terms is strictly less
than the leading term, so cancellation is impossible.
Choose a larger integer radius and subdivide rationally, preserving the
exact root-count sum.  If a proposed midpoint is a root, choose a nonroot
among the $r+1$ interior fractions $j/(r+2)$ of the interval; a degree-$r$
polynomial cannot vanish at all of them.  Every such split contracts each
child by at most $(r+1)/(r+2)$.  Distinct roots have positive separation,
so repeated contraction first isolates each and then attains the required
interval width.  This proves finite termination, including rational roots.

Evaluate $|A\phi|$ with an outward Arb ball covering both exact endpoints
of each isolating interval.  The largest upper endpoint is a global upper
certificate, because all maximising roots are covered.  Evaluation at
midpoints gives lower bounds on the global supremum even when those points
are not roots.  The independent implementation reconstructs the rational
polynomials, verifies total root counts and interval endpoint signs, and
uses higher-precision Gaussian evaluations to check the retained envelopes.
No claim that interval midpoint evaluation is the exact maximum is needed.

\section{Signed interpolation and arithmetic contracts}
\label{app:arithmetic}

Let $f\in C^2([a,a+s])$ and grid spacing $h=s/2^b$, with a power-of-two
span $s$ and $b\leq q+\log_2s$.  At every endpoint $x_j$, retain an
unsigned dyadic value $n_j2^{-q}$ and an outward error $e_j$.
On a cell $[u,v]$, the exact linear interpolation error equals
$f''(\zeta)(x-u)(x-v)/2$ for some interior $\zeta$.
Subtract the linear interpolant and a multiple of $(t-u)(t-v)$ chosen to
give a third zero at $x$, then apply Rolle's theorem twice to prove this
identity.  Since $(x-u)(v-x)\leq h^2/4$, endpoint rounding and one integer
truncation give
\begin{equation}
 e_f\leq\max_je_j+\|f''\|_\infty h^2/8+2^{-q}.
 \label{eq:appendix-interpolation}
\end{equation}
An argument error $\eta_x$ adds $\|f'\|_\infty\eta_x$ if both arguments
remain in the declared interval.  Projection onto that interval does not
enlarge the distance from an admissible exact argument.

The actual table stores the base integer, the absolute difference to the
next endpoint, and its decreasing bit.  It multiplies the unsigned magnitude
by the within-cell remainder and chooses unsigned addition or subtraction.
The adjustment cannot underflow a decreasing base: it is at most the
endpoint difference, and both endpoints are nonnegative.
The final endpoint has its own address, so saturation at the right endpoint
does not wrap to the left.  The packed magnitudes, intermediate additions
and input slices are checked against their actual widths.
Outward endpoints are coarsened downwards and upwards to
$2^{-384}\mathbb Z$ before serialisation; the enlarged intervals determine
all nodal errors.  Extreme Gaussian tails are not treated as exact zeros.

For $x\geq a>0$, usable first and second derivative bounds are
$a^{-1},a^{-2}$ for $\log x$, and $2a^{-3},6a^{-4}$ for $x^{-2}$.
For $e^x$ both are $e^{a+s}$.
For the boundary scalar,
$C_2'=C_3+2C_2$ and $C_2''=C_4+5C_3+4C_2$ give valid slope and
curvature bounds, divided by $2M_2$ for the capped $h$ kernel.
I intersect them with independent two-density derivative bounds when
smaller. Capped price tables receive $|P'|\leq M_1$,
$|P''|\leq M_1+M_2$ and $|P'''|\leq M_1+3M_2+M_3$, intersected with
$H\lambda^{-j}\|\phi^{(j-1)}\|_\infty$.
Vanilla raw-price tables use the finite-domain forward bound for $P'$;
its boundary-density scalar remains globally bounded. These inequalities
are shared across methods and never change a probability's decoder.

For products, either identity
$\widehat x\widehat y-xy=x(\widehat y-y)+\widehat y(\widehat x-x)
=y(\widehat x-x)+\widehat x(\widehat y-y)$ yields a valid absolute error.
Taking the smaller range-based upper bound and adding one truncation unit
preserves small-factor suppression.  Addition and subtraction add operand
errors.  Scaling charges its rounded coefficient and any dropped bits.
Each node independently propagates implemented ranges; intersecting them
with an ideal range expanded by its certified error retains correlations
without removing rounding.  Empty intersections or domain violations reject
the graph.  Signed final differences must fit strictly within the
unambiguous two's-complement range before projection.  Positive-subtraction
temporaries are predicate-masked before any downstream multiplication.
These checks are essential hypotheses of the monetary certificate.

\section{Auxiliary midpoint and stable quadrature proofs}
\label{app:aux}

For a Beta cell centred at $z_j$ with width $h$, write $u=z-z_j$.
The density is $2(1-z_j)-2u$; direct integration gives first centred
moment $-h^3/6$ and second centred moment $2(1-z_j)h^3/12$.
Taylor's formula, summed over cells with their exact positive masses,
therefore gives
\begin{equation}
 |\E f(Z_B)-\E f(Z_{B,N})|
 \leq h^2\{\|f'\|_\infty/6+\|f''\|_\infty/24\}.
\end{equation}
For $f_y(z)=g_{m(y)}''(1+zU)$, the derivatives are bounded by
$M_3U$ and $M_4U^2$. Multiplication by $U^2/2$ and finite financial
expectation prove \eqref{eq:beta-error}.

For the stable rule the Beta moments are $2/((j+1)(j+2))$.
The monic polynomial $\pi_2(z)=z^2-4z/5+1/10$ is orthogonal to
$1,z$, by substitution of these moments. Its roots and positive weights
are those in \eqref{eq:stable-kernel}; the first four moments show
exactness through degree three. Direct integration gives
$\int\pi_2^2\,2(1-z)dz=1/300$.
For $f\in C^4[0,1]$, its cubic Hermite interpolant $H_f$ agrees with
$f,f'$ at both nodes. At any non-node $z$, subtract a multiple of
$\pi_2^2$ chosen to create an additional zero there. Four applications of
Rolle's theorem, counting the two double zeros, give
$f(z)-H_f(z)=f^{(4)}(\xi_z)\pi_2(z)^2/4!$.
At a node the error is zero. Positivity and exactness of the quadrature
then imply error at most $\|f^{(4)}\|_\infty/7200$.
Taking $f=f_y$ bounds its fourth derivative by $M_6U^4$.
The factor $U^2/2$ gives the first term of \eqref{eq:gauss-error}.
Replace nodes first under the true positive weights, using $M_3U$;
then replace weights at the dyadic nodes, using $M_2$. This gives the two
remaining terms after finite expectation. All nodes lie in $[0,1]$.
The graph's dyadic constants are exact rational operations; their discrepancy
from the ideal quadrature is charged here once.

\section{Gaussian transport and factorised financial error}
\label{app:financial-transfer}

I complete the financial transfer for the four- and eight-asset models
by integrating derivative envelopes against a positive transport measure.
This retains dispersion moments instead of replacing $U$ by its largest
value on $[-L,L]^p$. Let $\nu$ be the unnormalised truncated Gaussian
measure, with mass $a$, and let $\tau=(a/N)\sum_c\delta_{c_c}$ use the
actual rounded nodes in the same CDF cells. For their centred cell moments
$\mu_{1,c},\mu_{2,c}$, the weighted identity of \citet{cerone2002}
(Theorem~3, $n=1$, $\alpha=a_c$, $\beta=b_c$, $x=c_c$, $w=\phi$) gives
\begin{equation}
 \int f\,d(\nu-\tau)=\sum_c\mu_{1,c}f'(c_c)
                         +\int f''\,d\rho_2,
\end{equation}
where $\rho_2$ has nonnegative densities
$\int_{a_c}^{z}(z-y)\phi(y)dy$ to the left of $c_c$ and
$\int_z^{b_c}(y-z)\phi(y)dy$ to its right. Its mass is
$\sum_c\mu_{2,c}/2$. For
$\rho_1=\sum_c\widehat\mu_{1,c}\delta_{c_c}$ with certified
$\widehat\mu_{1,c}\geq|\mu_{1,c}|$, the absolute transfer is bounded
by $\int|f'|d\rho_1+\int|f''|d\rho_2$.

For the two-asset local variant, I keep the signed linear sum and
bound the curvature on retained cell segments:
\begin{equation}
 \beta_Y\leq B_{\rm tail}+\left|\sum_c\mu_{1,c}R'(c_c)\right|
                    +\sum_{c,I}D_{c,I}\rho_2(I),
 \label{eq:reference-transfer}
\end{equation}
where $D_{c,I}$ encloses $\sup_I|R''|$. Nonnegative transport weights
justify outward endpoint extensions. The full local proof retains the
closed Gaussian residual derivatives, primitive weight integrals and
independent automatic-differentiation checks.

The exact moment-generating function of the second transport measure is
\begin{equation}
 \mathcal M_{\rho_2}(v)=\frac{
 e^{v^2/2}\{\Phi(L-v)-\Phi(-L-v)\}
 -(a/N)\sum_c e^{vc_c}-v\sum_c\mu_{1,c}e^{vc_c}}{v^2},
 \qquad v\ne0.
 \label{eq:positive-transport-mgf}
\end{equation}
This follows by substituting $f(y)=e^{vy}$ into the exact identity,
with continuous minus rounded measures in that order. At zero I use its
mass directly. Near zero, the enclosure
$e^{-|v|L}\rho_2(\mathbb R)\leq\mathcal M_{\rho_2}(v)
\leq e^{|v|L}\rho_2(\mathbb R)$ prevents cancellation in the quotient.

Define
$\Lambda_{j,k}=\nu^{\otimes(j-1)}\otimes\rho_k
\otimes\tau^{\otimes(p-j)}$ and
$\langle h\rangle_{j,k}=\int h\,d\Lambda_{j,k}$.
The product-measure telescoping identity then proves
\begin{equation}
 \left|\int R_2d\nu^{\otimes p}-\int R_2d\tau^{\otimes p}\right|
 \leq\sum_j\{\langle|\partial_jR_2|\rangle_{j,1}
                   +\langle|\partial_{jj}R_2|\rangle_{j,2}\}.
 \label{eq:positive-tensor-transfer}
\end{equation}
All raw masses remain inside the hybrid measures.

For $u=T_j$, $v=T_{jj}\geq0$ and $r=r_j$, let
$D_1\geq\|C_2'\|_\infty$, $D_2\geq\|C_2''\|_\infty$ and
$D_{31}\geq\|C_3'\|_\infty$, in addition to the shared $M_k$ bounds.
The integral representation of $R_2$ implies
\begin{align}
 |\partial_jR_2|&\leq M_2U|u|+\tfrac12|r|D_1U^2
                                  +\tfrac16M_3U^2|u|,\\
 |\partial_{jj}R_2|&\leq M_2(u^2+Uv)+2|r|D_1U|u|
 +\tfrac23M_3Uu^2+\tfrac12r^2D_2U^2\\
 &\quad+\tfrac13|r|D_{31}U^2|u|
                +\tfrac16M_3U^2v+\tfrac1{12}M_4U^2u^2.
\end{align}
Let $E_{j,1}$ and $E_{j,2}$ denote these two nonnegative derivative
envelopes. Under the stated positive-weight exponential geometry and
bounded curvature derivatives, they give the following full certificate.
\begin{theorem}[Factorised Gaussian financial transfer]
\label{thm:tensor-transfer}
For actual nodes certified inside their Gaussian CDF cells,
\begin{equation}
 |\E R_2-\alpha_1^p\E_NR_2|
 \leq \frac{M_2}{2}\E[U^2\ind_{Y\notin[-L,L]^p}]
 +\sum_j\{\langle E_{j,1}\rangle_{j,1}
                    +\langle E_{j,2}\rangle_{j,2}\}.
\end{equation}
All envelope terms are bounded by finite positive-measure moments of
exponential polynomials of degree at most four, using the Cauchy bounds
below. Their factorisation preserves the raw Gaussian masses.
\end{theorem}
\begin{proof}
Split the full Gaussian integral into the box and its complement.
The global curvature bound controls the latter by the displayed outside
moment. The positive one-coordinate representation, applied to each term
of the tensor telescope, controls the former by the two derivative
envelopes. Their derivation and finite-moment evaluation follow directly.
To verify the coefficients, differentiate
$G=t^{-2}C_2(m+\log t)$ with $t=1+zU$:
\begin{equation}
 G_j=rt^{-2}C_2'+zu t^{-3}C_3,
\quad G_{jj}=r^2t^{-2}C_2''+2rzu t^{-3}C_3'
                         +zv t^{-3}C_3+z^2u^2t^{-4}C_4.
\end{equation}
Apply the product rule to $U^2G$, use $t\geq1$, and integrate
$1-z$, $z(1-z)$ and $z^2(1-z)$, whose integrals are $1/2$, $1/6$
and $1/12$. Cauchy's inequality gives
$\langle U|u|\rangle\leq\sqrt{\langle U^2\rangle\langle u^2\rangle}$
and
\begin{equation}
 \langle U^2|u|\rangle\leq
 \min\{\sqrt{\langle U^2\rangle\langle U^2u^2\rangle},
          \sqrt{\langle U^4\rangle\langle u^2\rangle}\}.
\end{equation}
These inequalities hold for finite positive measures without a mass
divisor. Every remaining moment is a polynomial of degree at most four
in the asset exponentials. Its signed multinomial sum factorises into
scalar MGFs, using $\rho_k$ in coordinate $j$, $\nu$ before it and
$\tau$ after it. Thus no $N^p$ residual values enter the certificate.
The outside-box $U^2$ moment follows by Gaussian tilting and contributes
$M_2\mathbb E[U^2\mathbf1_{Y\notin[-L,L]^p}]/2$.

\end{proof}

For the matched FTC comparator, compact-cell Fubini applied to
$f(y)-f(c)$ gives a signed kernel $\Phi(a_c)-\Phi(z)$ on the left and
$\Phi(b_c)-\Phi(z)$ on the right. Its absolute density defines $\rho_0$.
Writing $J_v(a,b)=e^{v^2/2}\{\Phi(b-v)-\Phi(a-v)\}$, integration by parts
on both segments gives, for $v\ne0$,
\begin{equation}
 \mathcal M_{\rho_{0,c}}(v)=
 \frac{e^{vc}\{2\Phi(c)-\Phi(a_c)-\Phi(b_c)\}
             +J_v(c,b_c)-J_v(a_c,c)}{v}.
\end{equation}
At zero I use $\int_{a_c}^{b_c}|y-c|\phi(y)dy$ directly. Replacing
$\rho_1,\rho_2$ by $\rho_0$ and integrating $E_{j,1}$ gives its bound
for the same $R_2$. This independently derived common-grid adaptation
does not assert equivalence to the uniform-grid theorem of
\citet{herman2026}.

For any transported-coordinate block partition and nonnegative moment
pairs $(A_D,B_D)$, discrete Cauchy gives
$\sum_D\sqrt{A_DB_D}\leq\sqrt{(\sum_DA_D)(\sum_DB_D)}$.
The minimum of the two alternative Cauchy closures also cannot increase
under exact refinement, since $\sum_D\min(a_D,b_D)\leq
\min(\sum_Da_D,\sum_Db_D)$. All envelope coefficients are nonnegative
and direct moments add. Independently rounded numerical bounds are made
monotone by taking the minimum of parent and summed-child upper bounds;
both are valid. For $M$ monomials, the implemented shared grouping uses
$O(p^2M)$ other-coordinate products and $O(BpM)$ subsequent block dot
products, with scalar transforms, rate formation and prefix preparation
charged separately. Full cell transforms, interval-series remainders,
refinement proofs and independent audits are retained in the supporting
comparison record.

For a fair first-correction comparison, I also transfer
$R_1=R_2+UP'(m)$. The extra derivatives are
$uP'+rUP''$ and $vP'+2ruP''+r^2UP'''$. Capped kernels use the shared
global Greek bounds; vanilla kernels use
$P'\leq Ke^{m+\lambda^2/2}$, $|P''|\leq Ke^{m+\lambda^2/2}+M_2$
and $|P'''|\leq Ke^{m+\lambda^2/2}+3M_2+M_3$.
The corresponding tilted positive-measure moments also factorise.
Its tail is bounded by the tilted outside moment of
$Ke^{m+\lambda^2/2}U$, with the additional $H$ tail bound for a cap.
I replace the second-residual tail by this first-residual tail.

At unchanged $L=6$, $b_y=10$, $q=28$, all 24 four- and eight-asset
second-correction transfer bounds pass $0.005$; their largest allowance
is $0.002284$ rounded upwards. All 24 first-correction competitors
also pass; their largest allowance is $0.004914$ rounded upwards.
Each global-envelope family retains three two-asset allocation failures,
for which the earlier sharper local certificate remains available.
The transport argument applies to arbitrary finite positive basket
weights and bounded boundary derivatives, beyond these frozen models.
Its numerical scope and polynomial moment-construction cost do not
establish a quantum advantage or a literature-priority claim.

\section{Direct capped-price transfer and full confidence}
\label{app:direct-capped-price}

For the same positive-weight model let $f(y)=P(m_A(y))$ with
$m_A(y)=\log\sum_iw_i e^{\mu_i+R_i\cdot y}$ and $0\leq f\leq H$.
Define the positive weights $\pi_i=w_i e^{\mu_i+R_i\cdot y}/
\sum_kw_k e^{\mu_k+R_k\cdot y}$ and $A_j^*=\max_i|R_{ij}|$.
Differentiation gives $a_j=\partial_jm_A=\sum_i\pi_iR_{ij}$ and
$b_j=\partial_{jj}m_A=\sum_i\pi_iR_{ij}^2-a_j^2\geq0$. Combining
terms with $C_2=P''-P'$ yields
\begin{equation}
 f_{jj}=P'(m_A)(a_j^2+b_j)+C_2(m_A)a_j^2,
 \qquad a_j^2+b_j=\sum_i\pi_iR_{ij}^2\leq(A_j^*)^2.
\end{equation}
The shared global capped-kernel norms consequently imply
$|f_j|\leq M_1A_j^*$ and
$|f_{jj}|\leq(M_1+M_2)(A_j^*)^2$ on all real coordinates.
These use full asset loadings, not centred residual loadings.

For retained measure mass $a$, positive cell-error masses
$b_k=\rho_k(\R)$ and $F=a^p$, tensor telescoping proves
\begin{equation}
 \beta_{\rm cell}=a^{p-1}\sum_j
 \{M_1A_j^*b_1+(M_1+M_2)(A_j^*)^2b_2\},\quad
 V=F\mu_N+e_{\rm cell}+e_{\rm tail},
\end{equation}
where $|e_{\rm cell}|\leq\beta_{\rm cell}$,
$0\leq e_{\rm tail}\leq H(1-F)$, and $\mu_N$ is the direct
uniform-label mean at the actual rounded nodes. The global bounds grow
quadratically with exposure scale and do not require residual ratio moments.
They are a financial specialisation of established cell and derivative tools.

For independent bounded scores, \citet{hoeffding1963}, Theorem~2,
equation~(2.6), and the two-sided union in equation~(1.4), give
$\Pr\{|\overline f-\mu_N|>\epsilon\}\leq2e^{-2n\epsilon^2/H^2}$.
The full supporting proof derives this from the tilted variance for both
signs. I fix $n=100000$, $H=1$, $\epsilon=1/150$ before new draws.
Outward evaluation verifies that the failure bound is at most $1/3600$.
Masked 16-bit independent uniform words give exact 10-bit labels, each
with 64 preimages. Retained words verify this implementation; they do not
prove physical source independence.

Validated 128-bit scores are stored as outward 40-bit integer endpoints.
Exact summation gives $\overline\ell\leq\overline f\leq\overline u$.
On the confidence event, put $L_N=\max(0,\overline\ell-\epsilon)$ and
$U_N=\min(H,\overline u+\epsilon)$. With outward mass endpoints $F_-,F_+$,
\begin{equation}
 V\in[\max(0,F_-L_N-\widehat\beta_{\rm cell}),\;
 \min(H,F_+U_N+\widehat\beta_{\rm cell}+H(1-F_-))].
\end{equation}
All financial, numerical, statistical and mass errors enter this interval;
discounting is applied once afterwards. The exact finite sum replaces the
sample interval and removes the statistical term for the two-asset case.

The same bounds strengthen the first-correction opponent. For the normalised
weights, define $m_G(y)=\sum_iw_i(\mu_i+R_i\cdot y)=m(y)$ and
$R_{Gj}=\sum_iw_iR_{ij}$. Jensen and monotonicity give
$h=f-P(m_G)\in[0,H]$. Derivative subtraction bounds
$|h_j|$ by $M_1(A_j^*+|R_{Gj}|)$ and $|h_{jj}|$ by
$(M_1+M_2)((A_j^*)^2+R_{Gj}^2)$. The same cell formula and bounded tail
therefore certify $R_1$. Taking the minimum with its earlier valid transfer
is deterministic. The oracle and scalar words remain unchanged; monetary
and confidence plans use this stronger shared information. Full proofs,
all observed enclosures and independent replays are indexed separately.

\section{Global payoff-specific norm certificates}
\label{app:payoff-norms}

For the cap put $a=\ell/\lambda$ and
$q(z)=\phi(z)-c\phi(z-a)$. On $[0,a]$,
$q'=-z\phi(z)+c(z-a)\phi(z-a)<0$, including the endpoints.
Outside that interval, its stationary equation is equivalent to
\begin{equation}
 f(z)=\log(z/(z-a))-az-\ell+a^2/2=0.
\end{equation}
The ratio is positive on each component and no sign is lost in taking
logarithms. Its derivative is $-a/[z(z-a)]-a<-a$.
The limits on $(-\infty,0)$ are $+\infty,-\infty$, and on
$(a,\infty)$ they are again $+\infty,-\infty$. There is exactly one
root on each component. At a root,
$q(z)=-ca\phi(z-a)/z$, positive on the negative component and negative
on the positive one. Gaussian tails vanish. Thus
\begin{equation}
 \|C_2\|_\infty=\frac K\lambda
 \max\{q(z_-),-q(z_+)\}.
\end{equation}
Outward rational brackets with strict endpoint signs and the monotonicity
margin enclose both roots and their full stationary-value intervals.
An enclosed stationary midpoint gives a lower norm bound; the larger
stationary upper enclosure gives the required upper bound. Uncertainty in
$\ell$ and $a$ is retained throughout, rather than rounding them first.

For the capped first derivative,
$P'(m)=(K/\lambda)\int_0^\ell e^x\phi((m-x)/\lambda)dx>0$.
Normalising this integrand defines a truncated normal $X$ of mean parameter
$m+\lambda^2$ and variance parameter $\lambda^2$. Differentiation gives
\begin{equation}
 (\log P')'=(\E_mX-m)/\lambda^2,\qquad
 (\log P')''=-1/\lambda^2+\Var_m(X)/\lambda^4.
\end{equation}
For its density $p_m$ and mean $\bar x$, integration by parts yields
\begin{equation}
 -\Var_m(X)/\lambda^2
 =(\ell-\bar x)p_m(\ell)+\bar x p_m(0)-1.
\end{equation}
Both boundary terms are positive, since $0<\bar x<\ell$.
Therefore $\Var_m(X)<\lambda^2$ and $P'$ is strictly log-concave.
Its logarithmic derivative is positive at zero and negative at $\ell$,
so it has exactly one global maximum. Strict outward signs enclose that
zero of $P''$; evaluation on the whole bracket certifies $M_1$.
The supporting proof retains an explicit positive monotonicity margin.

For $r\geq2$, define
$\mathcal P_2=1$ and $\mathcal P_{r+1}=\lambda^{-1}(\mathcal P_r'-z\mathcal P_r)-r\mathcal P_r$.
The chain rule proves
$C_r^{\rm vanilla}(m)=K\mathcal P_r(m/\lambda)\phi(m/\lambda)/\lambda$;
the cap subtracts its $c$-weighted translate by $a$.
Exact stationary-root isolation for $\mathcal P_r'-z\mathcal P_r$ gives
$N_r=\|\mathcal P_r\phi\|_\infty$ outwards. Valid higher cap bounds are the
smaller of $H\overline D_r$ and $(2K+H)\overline N_r/\lambda$.
For $r=2$ the two-root result replaces this triangle bound. All competitors
receive the same constants, and the globally unbounded vanilla $P'$ is
bounded only on its certified finite argument domain.

\begin{thebibliography}{99}
\bibitem[Babbush et al.(2018)]{babbush} Babbush, R., Gidney, C., Berry, D. W., Wiebe, N., McClean, J., Paler, A., Fowler, A. and Neven, H. (2018). Encoding electronic spectra in quantum circuits with linear T complexity. \emph{Physical Review X}, 8(4), 041015. \url{https://doi.org/10.1103/PhysRevX.8.041015}.

\bibitem[Bayer et al.(2018)]{bayer} Bayer, C., Siebenmorgen, M. and Tempone, R. (2018). Smoothing the payoff for efficient computation of basket option prices. \emph{Quantitative Finance}, 18(3), 491--505. \url{https://doi.org/10.1080/14697688.2017.1308003}.

\bibitem[Blair et al.(2016)]{blair2016} Blair, J., Luttman, A. and Machorro, E. (2016). A generalized Peano Kernel Theorem for distributions of exponential decay. \emph{Journal of Mathematical Analysis and Applications}, 433(1), 622--641. \url{https://doi.org/10.1016/j.jmaa.2015.07.072}. Publisher abstract consulted; full-text access was unavailable.

\bibitem[Brassard et al.(2002)]{brassard} Brassard, G., H\o yer, P., Mosca, M. and Tapp, A. (2002). Quantum amplitude amplification and estimation. \emph{Contemporary Mathematics}, 305, 53--74. \url{https://doi.org/10.1090/conm/305/05215}.

\bibitem[Caldana et al.(2016)]{caldana}
Caldana, R., Fusai, G., Gnoatto, A. and Grasselli, M. (2016).
General closed-form basket option pricing bounds.
\emph{Quantitative Finance}, 16(4), 535--554.
\url{https://doi.org/10.1080/14697688.2015.1073854}. Detailed version consulted: City Research Online accepted manuscript, Sections 1--3 and Appendices B and D.

\bibitem[Carr(2001)]{carr2001} Carr, P. (2001). Deriving derivatives of derivative securities. \emph{Journal of Computational Finance}, 4(2), Winter 2000/2001. \url{https://doi.org/10.21314/JCF.2001.075}.

\bibitem[Carrera Vazquez and Woerner(2021)]{carreravazquez2021}
Carrera Vazquez, A. and Woerner, S. (2021). Efficient state preparation for quantum amplitude estimation.
\emph{Physical Review Applied}, 15, 034027. \url{https://doi.org/10.1103/PhysRevApplied.15.034027}.
Detailed version consulted: arXiv:2005.07711v1, 15 May 2020.

\bibitem[Cerone(2002)]{cerone2002} Cerone, P. (2002). Weighted Three Point Identities and their Bounds. \emph{SUT Journal of Mathematics}, 38(1), 17--37. \url{https://www.rs.tus.ac.jp/sutjmath/_userdata/38-1/02-cerone.pdf}. Final published PDF consulted, Theorems 3 and 5, pages 23--30.

\bibitem[Chakrabarti et al.(2021)]{chakrabarti} Chakrabarti, S., Krishnakumar, R., Mazzola, G., Stamatopoulos, N., Woerner, S. and Zeng, W. J. (2021). A threshold for quantum advantage in derivative pricing. \emph{Quantum}, 5, 463. \url{https://doi.org/10.22331/q-2021-06-01-463}.

\bibitem[Chen et al.(2026)]{chenlineufeld2026} Chen, J., Li, Y. and Neufeld, A. (2026). Quantum Monte Carlo algorithm for option pricing and its complexity analysis. \emph{Mathematical Finance}, Early View, first published 10 June 2026. \url{https://doi.org/10.1111/mafi.70036}. Detailed version consulted: arXiv:2301.09241v6, 3 May 2026.

\bibitem[Curran(1994)]{curran} Curran, M. (1994). Valuing Asian and portfolio options by conditioning on the geometric mean price. \emph{Management Science}, 40(12), 1705--1711. \url{https://doi.org/10.1287/mnsc.40.12.1705}. Final journal scan, pages 1705--1711, obtained from the author upload through Safari.

\bibitem[Dinge\c{c} and H\"ormann(2013)]{dingechormann} Dinge\c{c}, K. D. and H\"ormann, W. (2013). Control variates and conditional Monte Carlo for basket and Asian options. \emph{Insurance: Mathematics and Economics}, 52(3), 421--434. \url{https://doi.org/10.1016/j.insmatheco.2013.03.002}.

\bibitem[Duchi et al.(2012)]{duchi2012}
Duchi, J. C., Bartlett, P. L. and Wainwright, M. J. (2012). Randomized smoothing for stochastic optimization.
\emph{SIAM Journal on Optimization}, 22(2), 674--701. \url{https://doi.org/10.1137/110831659}.
Detailed version consulted: arXiv:1103.4296v2, 7 April 2012.

\bibitem[Glasserman and Yao(1992)]{glassermanyao1992}
Glasserman, P. and Yao, D. D. (1992). Some guidelines and guarantees for common random numbers.
\emph{Management Science}, 38(6), 884--908. \url{https://doi.org/10.1287/mnsc.38.6.884}.

\bibitem[Herman et al.(2026)]{herman2026} Herman, D., Sun, Y., Liu, J.-P., Pistoia, M., Che, C., Otter, R., Chakrabarti, S. and Harrow, A. (2026). Quantum Speedups for Derivative Pricing Beyond Black-Scholes. Unpublished preprint, arXiv:2602.03725v1, 3 February. \url{https://arxiv.org/html/2602.03725v1}. No published successor identified in the retained source audit.

\bibitem[Hoeffding(1963)]{hoeffding1963} Hoeffding, W. (1963). Probability inequalities for sums of bounded random variables. \emph{Journal of the American Statistical Association}, 58(301), 13--30. \url{https://doi.org/10.1080/01621459.1963.10500830}. Final journal scan consulted, printed pages 13--16, particularly Theorem 2, equation (2.6).

\bibitem[Imai and Tan(2006)]{imai2006}
Imai, J. and Tan, K. S. (2006). A general dimension reduction technique for
derivative pricing. \emph{Journal of Computational Finance}, 10(2), 129--155.
\url{https://doi.org/10.21314/JCF.2006.143}. Material consulted: final publisher abstract and author-uploaded final journal page 129.

\bibitem[Johansson(2017)]{johansson2017} Johansson, F. (2017). Arb: efficient arbitrary-precision midpoint-radius interval arithmetic. \emph{IEEE Transactions on Computers}, 66(8), 1281--1292. \url{https://doi.org/10.1109/TC.2017.2690633}.

\bibitem[Ju(2002)]{ju}
Ju, N. (2002).
Pricing Asian and basket options via Taylor expansion.
\emph{Journal of Computational Finance}, 5(3), 79--103.
\url{https://doi.org/10.21314/JCF.2002.088}.

\bibitem[Marinelli(2024)]{marinelli2024}
Marinelli, C. (2024). On certain representations of pricing functionals.
\emph{Annals of Finance}, 20, 91--127. \url{https://doi.org/10.1007/s10436-024-00438-5}.

\bibitem[Montanaro(2015)]{montanaro} Montanaro, A. (2015). Quantum speedup of Monte Carlo methods. \emph{Proceedings of the Royal Society A}, 471(2181), 20150301. \url{https://doi.org/10.1098/rspa.2015.0301}.

\bibitem[Olivares and Alvarez(2016)]{olivaresalvarez2016} Olivares, P. and Alvarez, A. (2016). Pricing basket options by polynomial approximations. \emph{Journal of Applied Mathematics}, 2016, 9747394. \url{https://doi.org/10.1155/2016/9747394}.

\bibitem[Rall and Fuller(2023)]{rallfuller2023} Rall, P. and Fuller, B. (2023). Amplitude estimation from quantum signal processing. \emph{Quantum}, 7, 937. \url{https://doi.org/10.22331/q-2023-03-02-937}.

\bibitem[Rebentrost et al.(2018)]{rebentrost} Rebentrost, P., Gupt, B. and Bromley, T. R. (2018). Quantum computational finance: Monte Carlo pricing of financial derivatives. \emph{Physical Review A}, 98, 022321. \url{https://doi.org/10.1103/PhysRevA.98.022321}.

\bibitem[Ross and Selinger(2016)]{rossselinger2016}
Ross, N.J. and Selinger, P. (2016).
Optimal ancilla-free Clifford+T approximation of z-rotations.
\emph{Quantum Information and Computation}, 16(11--12), 901--953.
\url{https://doi.org/10.26421/QIC16.11-12-1}.

\bibitem[Stamatopoulos et al.(2020)]{stamatopoulos2020} Stamatopoulos, N., Egger, D. J., Sun, Y., Zoufal, C., Iten, R., Shen, N. and Woerner, S. (2020). Option pricing using quantum computers. \emph{Quantum}, 4, 291. \url{https://doi.org/10.22331/q-2020-07-06-291}.

\bibitem[Stamatopoulos and Zeng(2024)]{stamatopoulos} Stamatopoulos, N. and Zeng, W. J. (2024). Derivative pricing using quantum signal processing. \emph{Quantum}, 8, 1322. \url{https://doi.org/10.22331/q-2024-04-30-1322}.

\bibitem[Suzuki et al.(2020)]{suzuki2020}
Suzuki, Y., Uno, S., Raymond, R., Tanaka, T., Onodera, T. and Yamamoto, N. (2020).
Amplitude estimation without phase estimation. \emph{Quantum Information Processing}, 19, 75.
\url{https://doi.org/10.1007/s11128-019-2565-2}.
\end{thebibliography}
\end{document}
